\documentclass[11pt,a4paper,oneside]{article}

\usepackage[utf8]{inputenc}
\usepackage[T1]{fontenc}
\usepackage[spanish,es-tabla,es-nodecimaldot,es-noshorthands]{babel}
\spanishplainpercent
\usepackage{lmodern}
\usepackage{microtype}
\usepackage[a4paper,top=2.5cm,bottom=2.5cm,left=2.7cm,right=2.7cm,headheight=14pt]{geometry}
\usepackage{amsmath,amssymb,mathtools,bm}
\usepackage{booktabs,tabularx,multirow,array}
\usepackage{graphicx}
\usepackage[table]{xcolor}
\usepackage{tikz}
\usetikzlibrary{positioning,arrows.meta,calc,fit,backgrounds,shapes.geometric,shapes.misc,decorations.pathreplacing,babel}
\usepackage{pgfplots}
\pgfplotsset{compat=1.18}
\usepgfplotslibrary{groupplots}
\usepackage[ruled,vlined,linesnumbered,spanish]{algorithm2e}
\usepackage{caption}
\usepackage{enumitem}
\usepackage{fancyhdr}
\usepackage{csquotes}
\usepackage[backend=biber,style=numeric-comp,sorting=none,maxbibnames=6,giveninits=true,url=false,doi=true]{biblatex}
\addbibresource{referencias.bib}
\usepackage[colorlinks=true,linkcolor=azulenlace,citecolor=azulenlace,urlcolor=azulenlace]{hyperref}
\usepackage[spanish,noabbrev]{cleveref}

\definecolor{azulenlace}{HTML}{1F4E79}
\definecolor{cazul}{HTML}{0072B2}
\definecolor{cbermellon}{HTML}{D55E00}
\definecolor{cverde}{HTML}{009E73}
\definecolor{cnaranja}{HTML}{E69F00}
\definecolor{ccielo}{HTML}{56B4E9}
\definecolor{cpurpura}{HTML}{CC79A7}
\definecolor{cgris}{HTML}{6E6E6E}
\definecolor{fondogen}{HTML}{EAF3FA}
\definecolor{fondocat}{HTML}{FDF1E3}

\crefname{table}{tabla}{tablas}
\Crefname{table}{Tabla}{Tablas}
\crefname{figure}{figura}{figuras}
\Crefname{figure}{Figura}{Figuras}
\crefname{algocf}{algoritmo}{algoritmos}
\Crefname{algocf}{Algoritmo}{Algoritmos}
\crefname{equation}{ecuación}{ecuaciones}
\Crefname{equation}{Ecuación}{Ecuaciones}
\crefname{section}{sección}{secciones}
\Crefname{section}{Sección}{Secciones}
\crefname{subsection}{apartado}{apartados}
\Crefname{subsection}{Apartado}{Apartados}
\crefname{appendix}{apéndice}{apéndices}
\Crefname{appendix}{Apéndice}{Apéndices}

\captionsetup{font=small,labelfont=bf,format=hang,margin=0.5cm}
\setlist{itemsep=2pt,topsep=4pt}
\setcounter{tocdepth}{2}
\renewcommand{\topfraction}{0.85}
\renewcommand{\bottomfraction}{0.6}
\renewcommand{\textfraction}{0.12}
\renewcommand{\floatpagefraction}{0.75}
\makeatletter
\setlength{\@fptop}{0pt}
\setlength{\@fpsep}{14pt plus 2fil}
\setlength{\@fpbot}{0pt plus 1fil}
\makeatother
\renewcommand{\arraystretch}{1.15}
\SetAlCapSkip{0.6em}
\SetKwInput{KwIn}{Entrada}
\SetKwInput{KwOut}{Salida}
\SetKw{KwTo}{hasta}
\SetKw{KwRet}{devolver}
\SetKwFor{For}{para}{hacer}{fin}
\SetKwFor{ForEach}{para cada}{hacer}{fin}
\SetKwFor{While}{mientras}{hacer}{fin}
\SetKwIF{If}{ElseIf}{Else}{si}{entonces}{si no, si}{si no}{fin}

\pgfplotsset{
  informe/.style={
    width=\linewidth,
    height=6.2cm,
    grid=major,
    grid style={line width=0.2pt,draw=gray!25},
    axis line style={draw=gray!70},
    tick style={draw=gray!70},
    tick label style={font=\footnotesize},
    label style={font=\small},
    legend style={font=\footnotesize,draw=gray!50,fill=white,fill opacity=0.9,text opacity=1},
    legend cell align=left,
    every axis plot/.append style={line join=round},
  },
}

\DeclareMathOperator*{\argmin}{arg\,m\acute{\imath}n}
\DeclareMathOperator*{\argmax}{arg\,m\acute{a}x}
\DeclareMathOperator{\DiceOp}{Dice}
\DeclareMathOperator{\IoUOp}{IoU}
\newcommand{\R}{\mathbb{R}}
\newcommand{\E}{\mathbb{E}}
\newcommand{\vtheta}{\bm{\theta}}
\newcommand{\vlambda}{\bm{\lambda}}
\newcommand{\vu}{\mathbf{u}}
\newcommand{\vx}{\mathbf{x}}
\newcommand{\vy}{\mathbf{y}}
\newcommand{\vz}{\mathbf{z}}
\newcommand{\vp}{\mathbf{p}}
\newcommand{\vg}{\mathbf{g}}
\newcommand{\Loss}{\mathcal{L}}
\newcommand{\Dtrain}{\mathcal{D}_{\mathrm{ent}}}
\newcommand{\Dval}{\mathcal{D}_{\mathrm{val}}}
\newcommand{\Dtest}{\mathcal{D}_{\mathrm{test}}}
\newcommand{\aptitud}{\widehat{F}}

\newcommand{\universidad}{Universidad de Las Palmas de Gran Canaria}
\newcommand{\asignatura}{Optimización y Heurística}
\newcommand{\cursoacademico}{Curso 2026--2027}
\newcommand{\tituloinforme}{Optimización en dos niveles para la segmentación automática de tumores cerebrales en resonancia magnética}
\newcommand{\subtituloinforme}{Descenso de gradiente en una red U-Net y metaoptimización de hiperparámetros mediante un algoritmo genético}
\newcommand{\fechainforme}{Noviembre de 2026}

\hypersetup{pdftitle={\tituloinforme},pdfauthor={Javier Ruano Hernández, Alejandro Delgado Valera},pdfsubject={\asignatura}}

\pagestyle{fancy}
\fancyhf{}
\fancyhead[L]{\small\itshape \asignatura}
\fancyhead[R]{\small\itshape Segmentación de tumores cerebrales}
\fancyfoot[C]{\small\thepage}
\renewcommand{\headrulewidth}{0.4pt}

\begin{document}

\begin{titlepage}
  \centering
  \vspace*{1.2cm}
  {\large\scshape \universidad\par}
  \vspace{0.3cm}
  {\normalsize Asignatura: \asignatura\par}
  {\normalsize \cursoacademico\par}
  \vspace{2.6cm}
  \rule{\linewidth}{0.8pt}\par
  \vspace{0.5cm}
  {\LARGE\bfseries \tituloinforme\par}
  \vspace{0.5cm}
  {\large\itshape \subtituloinforme\par}
  \vspace{0.4cm}
  \rule{\linewidth}{0.8pt}\par
  \vspace{2.2cm}
  {\large Memoria final del proyecto\par}
  \vspace{1.6cm}
  \begin{tabular}{c}
    \large Javier Ruano Hernández\\[4pt]
    \large Alejandro Delgado Valera
  \end{tabular}\par
  \vfill
  {\large \fechainforme\par}
\end{titlepage}

\begin{abstract}
\noindent
La segmentación automática de tumores cerebrales en resonancia magnética es un problema de visión por computador con un fuerte desequilibrio de clases: en el conjunto de datos utilizado, formado por 3064 cortes ponderados en T1 con contraste de 233 pacientes, el tumor ocupa de media el 1.7\,\% de los píxeles de cada imagen. Este trabajo aborda la segmentación multiclase de meningiomas, gliomas y tumores hipofisarios como un \emph{problema de optimización en dos niveles}. En el nivel interno, el entrenamiento de una red U-Net se formula como la minimización de una función de pérdida diferenciable y no convexa sobre entre $1.8$ y $7.2$ millones de parámetros, resuelta mediante descenso de gradiente estocástico adaptativo (Adam) con gradientes exactos obtenidos por retropropagación. Se analizan desde el punto de vista de la optimización la arquitectura, las funciones de pérdida sustitutas (Dice, Focal y Focal Tversky) y la programación de la tasa de aprendizaje. En el nivel externo, los hiperparámetros que condicionan ese proceso (tasa de aprendizaje, decaimiento de pesos, número de filtros y profundidad) se ajustan con un algoritmo genético de codificación real, con selección por torneo, cruce BLX-$\alpha$, mutación gaussiana con reflexión y elitismo, y una función de aptitud de baja fidelidad basada en el coeficiente Dice de validación tras tres épocas de entrenamiento. Con 46 evaluaciones, el algoritmo genético mejora la aptitud de la configuración manual de 0.441 a 0.509. Reentrenada a longitud completa, la configuración evolucionada eleva el Dice medio en el conjunto de prueba de 0.691 a 0.776 y alcanza el mejor valor de validación de la referencia en 10 épocas en lugar de 35. El análisis posterior de la solución heurística revela un sesgo de horizonte corto en la función de aptitud, que favorece redes pequeñas de convergencia rápida. Corregido mediante recocido coseno de la tasa de aprendizaje, mayor capacidad (32 filtros) y una calibración de los pesos de la pérdida Focal Tversky guiada por diagnóstico, el modelo entrenado desde cero alcanza un Dice medio de 0.8175 en prueba. Como complemento, un codificador ResNet34 preentrenado en ImageNet obtiene 0.8250, un resultado estadísticamente equivalente que resuelve la confusión entre tipos tumorales pero no mejora la delimitación del glioma. El algoritmo genético explica aproximadamente dos tercios de la mejora total respecto a la configuración manual.

\medskip
\noindent\textbf{Palabras clave:} optimización en dos niveles, descenso de gradiente, retropropagación, U-Net, algoritmo genético, optimización de hiperparámetros, pérdida Focal Tversky, segmentación de tumores cerebrales.
\end{abstract}


\clearpage
\tableofcontents
\clearpage

\section{Introducción}\label{sec:introduccion}

\subsection{Contexto clínico}\label{subsec:contexto}

Los tumores del sistema nervioso central constituyen un grupo heterogéneo de neoplasias cuyo diagnóstico, planificación terapéutica y seguimiento dependen de forma crítica de la imagen médica~\cite{louis2021who}. La resonancia magnética (RM) es la modalidad de referencia para su estudio, por su elevado contraste entre tejidos blandos y por no emplear radiación ionizante. Entre las secuencias habituales, la ponderada en T1 tras la administración de un contraste de gadolinio (T1-CE) realza las regiones en las que la barrera hematoencefálica está alterada, lo que hace visible el tejido tumoral activo frente al parénquima sano. Cada corte de esta secuencia es una imagen bidimensional en escala de grises: un único canal de intensidad, sin información de color, en el que la lesión se distingue exclusivamente por su brillo, su textura y su forma.

Este trabajo considera tres tipos de tumor cerebral con características radiológicas muy distintas:

\begin{itemize}
  \item \textbf{Meningioma.} Tumor extraaxial, generalmente benigno, que se origina en las meninges. Suele presentar bordes bien definidos y un realce intenso y homogéneo, aunque puede aparecer en múltiples localizaciones de la convexidad y de la base del cráneo.
  \item \textbf{Glioma.} Tumor intraaxial originado en las células gliales del parénquima cerebral. Su crecimiento infiltrativo produce márgenes difusos e irregulares y un realce heterogéneo, a menudo con zonas necróticas. Es, con diferencia, el tipo más difícil de delimitar, también para los especialistas.
  \item \textbf{Tumor hipofisario.} Lesión localizada en la región selar, en la línea media de la base del cráneo. Suele ser de pequeño tamaño y de contornos regulares, pero su proximidad a otras estructuras de la base craneal puede confundirlo con un meningioma de esa zona.
\end{itemize}

La delimitación manual del tumor corte a corte es una tarea lenta y sujeta a variabilidad entre observadores. Un sistema de segmentación automática que proporcione una máscara precisa de la lesión y de su tipo permitiría cuantificar volúmenes de forma reproducible, acelerar la planificación quirúrgica o radioterápica y estandarizar el seguimiento longitudinal.

\subsection{Planteamiento del problema}\label{subsec:problema}

Formalmente, dado un corte $\vx \in [0,1]^{H\times W}$ normalizado, el objetivo es asignar a cada píxel $i$ una etiqueta $y_i \in \{0,1,2,3\}$, donde $0$ denota el fondo (incluido el tejido sano) y $1$, $2$ y $3$ denotan meningioma, glioma y tumor hipofisario, respectivamente. Se trata, por tanto, de un problema de \emph{segmentación semántica multiclase} en el que cada corte contiene un único tipo tumoral, de modo que la red debe resolver simultáneamente dónde está el tumor y de qué tipo es.

El problema presenta cuatro dificultades que condicionan todas las decisiones de diseño posteriores:

\begin{enumerate}
  \item \textbf{Desequilibrio extremo de clases.} En las imágenes originales, el tumor ocupa de media el 1.7\,\% del corte (mediana del 1.3\,\%, y solo el 0.8\,\% en los tumores hipofisarios). Un clasificador trivial que asignase todos los píxeles al fondo alcanzaría una exactitud superior al 98\,\% y, sin embargo, un coeficiente Dice nulo para las tres clases tumorales. La función objetivo debe, por tanto, diseñarse explícitamente para no quedar dominada por el fondo.
  \item \textbf{Variabilidad morfológica.} El tamaño, la forma, la localización y la intensidad de las lesiones varían mucho entre pacientes y entre los planos axial, coronal y sagital presentes en el conjunto de datos.
  \item \textbf{Similitud entre clases.} Meningiomas de la base del cráneo y tumores hipofisarios comparten localización y apariencia en un único corte bidimensional, y los bordes infiltrativos del glioma se confunden con el tejido circundante.
  \item \textbf{Escasez de datos.} Se dispone de 3064 cortes de 233 pacientes, un volumen reducido para entrenar desde cero una red convolucional profunda.
\end{enumerate}

\subsection{Enfoque: un problema de optimización en dos niveles}\label{subsec:enfoque}

La tesis central de este trabajo es que la construcción de un segmentador eficaz es, en su totalidad, un problema de optimización con dos niveles anidados, y que ambos niveles merecen el mismo tratamiento riguroso.

En el \textbf{nivel interno}, una vez fijado un vector de hiperparámetros $\vlambda$, el entrenamiento de la red consiste en encontrar los pesos $\vtheta$ que minimizan una función de pérdida diferenciable sobre el conjunto de entrenamiento. Es un problema de optimización continua de muy alta dimensión (entre $10^6$ y $10^7$ variables de decisión en este trabajo), no convexo y estocástico, que se resuelve con un método de primer orden cuyos gradientes se calculan de forma exacta mediante retropropagación. La arquitectura U-Net~\cite{ronneberger2015unet} no es aquí un mero modelo predictivo, sino la parametrización que define el espacio de búsqueda y la geometría del paisaje de la función objetivo.

En el \textbf{nivel externo}, los hiperparámetros $\vlambda$ que gobiernan el proceso anterior (la tasa de aprendizaje, la regularización y el tamaño de la red) se eligen para maximizar la calidad del modelo resultante sobre datos no vistos. Este segundo problema no admite gradientes, mezcla variables continuas y categóricas, y cada evaluación de la función objetivo exige entrenar una red completa. Es el escenario natural de las metaheurísticas, y en este trabajo se aborda con un algoritmo genético (AG)~\cite{holland1975adaptation,goldberg1989genetic}.

Ambos niveles se relacionan mediante la formulación binivel
\begin{equation}\label{eq:binivel}
  \vlambda^{\star} \in \argmax_{\vlambda \in \Lambda} \; \DiceOp_{\mathrm{val}}\bigl(\vtheta^{\star}(\vlambda)\bigr)
  \quad \text{sujeto a} \quad
  \vtheta^{\star}(\vlambda) \in \argmin_{\vtheta} \; \Loss_{\mathrm{ent}}(\vtheta;\vlambda),
\end{equation}
donde $\Lambda$ es el espacio de búsqueda de hiperparámetros, $\Loss_{\mathrm{ent}}$ la pérdida de entrenamiento y $\DiceOp_{\mathrm{val}}$ el coeficiente Dice medio sobre el conjunto de validación. El problema interno se resuelve de forma aproximada por descenso de gradiente, y el externo de forma aproximada por evolución artificial.

\subsection{Objetivos}\label{subsec:objetivos}

Los objetivos del proyecto son los siguientes:

\begin{enumerate}
  \item Formular e implementar el entrenamiento de una U-Net multiclase como un problema de optimización continua, con funciones de pérdida sustitutas diseñadas para el desequilibrio de clases (Dice-Focal y Focal Tversky), y analizar el papel de la arquitectura, el optimizador y la programación de la tasa de aprendizaje en ese proceso.
  \item Diseñar e implementar un algoritmo genético de codificación real que optimice los hiperparámetros de la red mediante una función de aptitud de bajo coste.
  \item Comparar analíticamente la configuración obtenida por la heurística con una configuración manual de referencia, tanto en precisión (Dice e IoU sobre un conjunto de prueba independiente) como en velocidad de convergencia.
  \item Analizar críticamente las limitaciones de la solución heurística y superarlas mediante experimentos controlados.
  \item De forma complementaria, evaluar si la transferencia de aprendizaje desde ImageNet supera los límites intrínsecos de un conjunto de datos pequeño, y ofrecer una interfaz gráfica para la inferencia interactiva.
\end{enumerate}

\subsection{Contribuciones y estructura del documento}\label{subsec:estructura}

Las principales aportaciones del trabajo son: (i) un análisis del entrenamiento de la U-Net desde la óptica de la optimización, que incluye el gradiente de la pérdida Focal Tversky y su efecto sobre el equilibrio entre precisión y sensibilidad; (ii) un algoritmo genético propio, con codificación real normalizada, control de novedad y memoización de la aptitud, que mejora la configuración manual en 0.086 puntos de Dice medio en prueba; (iii) la identificación de un sesgo de horizonte corto en la aptitud del AG y su corrección, que eleva el modelo entrenado desde cero hasta un Dice medio de 0.8175 en prueba; y (iv) una evaluación honesta de la transferencia de aprendizaje con intervalos de confianza por remuestreo de pacientes.

El documento se organiza como sigue. La \cref{sec:datos} describe el conjunto de datos, el preprocesamiento y el protocolo experimental. La \cref{sec:unet} desarrolla el nivel interno: la U-Net como problema de optimización resuelto por descenso de gradiente y retropropagación. La \cref{sec:ag} desarrolla el nivel externo: el algoritmo genético, su codificación, su función de aptitud, sus operadores y sus resultados. La \cref{sec:refinamiento} analiza las limitaciones de la solución heurística y describe cómo se superaron. La \cref{sec:transferencia} presenta, como complemento, el experimento de transferencia de aprendizaje. Por último, la \cref{sec:discusion} discute los resultados, sus amenazas a la validez y las líneas de trabajo futuro.

\section{Datos y protocolo experimental}\label{sec:datos}

\subsection{Conjunto de datos}\label{subsec:dataset}

Se emplea el conjunto público de Figshare~\cite{cheng2017figshare}, introducido por Cheng et~al.~\cite{cheng2015enhanced}, que contiene 3064 cortes de RM ponderada en T1 con contraste procedentes de 233 pacientes: 708 cortes de meningioma (82 pacientes), 1426 de glioma (89 pacientes) y 930 de tumor hipofisario (62 pacientes). Cada registro incluye la imagen, de 16 bits de profundidad y $512\times512$ píxeles (salvo 15 cortes de $256\times256$), el tipo de tumor, el identificador del paciente y una máscara binaria de la lesión delineada manualmente por radiólogos. Los cortes proceden de los planos axial, coronal y sagital. Para el problema multiclase, la máscara binaria se combina con el tipo de tumor, de modo que cada píxel tumoral recibe la etiqueta de su clase ($1$, $2$ o $3$) y el resto la etiqueta de fondo ($0$).

\subsection{Partición por paciente}\label{subsec:particion}

Cortes adyacentes de un mismo paciente son casi idénticos, por lo que una partición aleatoria por corte filtraría información entre subconjuntos e inflaría artificialmente las métricas. Por ello, todas las particiones se realizan \emph{por paciente} y estratificadas por tipo de tumor, con semilla fija. Primero se reserva el 20\,\% de los pacientes de cada clase como conjunto de prueba, y después el 20\,\% de los pacientes restantes como conjunto de validación. La \cref{tab:particion} resume el resultado.

\begin{table}[htbp]
  \centering
  \caption{Partición por paciente, estratificada por tipo de tumor. Entre paréntesis, pacientes de meningioma, glioma y tumor hipofisario.}
  \label{tab:particion}
  \begin{tabular}{lcrrrr}
    \toprule
    \textbf{Subconjunto} & \textbf{Pacientes} & \textbf{Meningioma} & \textbf{Glioma} & \textbf{Hipofisario} & \textbf{Total} \\
    \midrule
    Entrenamiento & 150 (53/57/40) & 461 & 938 & 605 & 2004 \\
    Validación    & 37 (13/14/10)  & 118 & 188 & 162 & 468 \\
    Prueba        & 46 (16/18/12)  & 129 & 300 & 163 & 592 \\
    \midrule
    Total         & 233 (82/89/62) & 708 & 1426 & 930 & 3064 \\
    \bottomrule
  \end{tabular}
\end{table}

El conjunto de prueba no interviene en el entrenamiento, en la selección de puntos de control ni en la búsqueda del algoritmo genético. El conjunto de validación cumple tres funciones: es el entorno en el que se mide la aptitud de cada individuo del AG, el criterio para seleccionar el mejor punto de control de cada entrenamiento y la base de los diagnósticos de error.

\subsection{Preprocesamiento}\label{subsec:preprocesamiento}

Cada par imagen-máscara atraviesa la siguiente secuencia determinista:

\begin{enumerate}
  \item \textbf{Recorte de márgenes.} Se detecta el contorno de la cabeza mediante un suavizado gaussiano, una umbralización al 8\,\% del rango de intensidad y una apertura morfológica que elimina ruido y anotaciones incrustadas. El recorte se calcula como la caja que engloba todas las componentes conexas relevantes y la propia lesión, con un margen de 4 píxeles. Así se elimina fondo negro que no aporta información y se aumenta la resolución efectiva del cerebro.
  \item \textbf{Redimensionado} a $256\times256$ píxeles, con interpolación por área para la imagen y por vecino más próximo para la máscara, de modo que cada píxel de la máscara conserva una etiqueta válida.
  \item \textbf{Normalización min-max} al intervalo $[0,1]$, tras recortar las intensidades a los percentiles 0.5 y 99.5 para que unos pocos píxeles saturados no compriman el rango útil del tejido.
\end{enumerate}

\subsection{Aumento de datos}\label{subsec:aumento}

El aumento de datos actúa como regularizador del problema de optimización interno: amplía el conjunto efectivo de entrenamiento con transformaciones que preservan la etiqueta. Se implementaron dos bloques independientes con la biblioteca Albumentations~\cite{buslaev2020albumentations}, combinados en un diseño factorial $2\times2$ (\cref{tab:aumento}):

\begin{itemize}
  \item \textbf{Bloque espacial básico}: volteos horizontal y vertical ($p=0.5$ cada uno) y rotación de hasta $\pm15^\circ$ en todas las muestras.
  \item \textbf{Bloque espacial avanzado}: volteo horizontal ($p=0.5$); transformación afín con rotación de $\pm15^\circ$, escala entre $0.9$ y $1.1$ y traslación de $\pm6\,\%$ ($p=0.7$); y, con $p=0.3$, una deformación elástica ($\alpha=200$, $\sigma=8$) o una distorsión de rejilla (5 pasos, límite $0.15$). Se suprime el volteo vertical porque el conjunto contiene cortes sagitales y coronales, en los que una inversión arriba-abajo genera anatomía inexistente.
  \item \textbf{Bloque de intensidad}: brillo y contraste ($\pm0.15$, $p=0.5$), corrección gamma (factor entre 0.8 y 1.2, $p=0.3$) y ruido gaussiano (desviación entre 0.01 y 0.03, $p=0.2$), aplicados tras la normalización.
\end{itemize}

\begin{table}[htbp]
  \centering
  \caption{Diseño factorial de los preajustes de aumento de datos. Cada preajuste de aislamiento difiere de los demás en exactamente un bloque.}
  \label{tab:aumento}
  \begin{tabular}{lcc}
    \toprule
     & \textbf{Sin intensidad} & \textbf{Con intensidad} \\
    \midrule
    \textbf{Espacial básico}   & \texttt{basic}            & \texttt{advanced\_intensity} \\
    \textbf{Espacial avanzado} & \texttt{advanced\_spatial} & \texttt{advanced} \\
    \bottomrule
  \end{tabular}
\end{table}

\subsection{Métricas de evaluación}\label{subsec:metricas}

Sean $P_c$ y $G_c$ los conjuntos de píxeles predichos y reales de la clase $c$, acumulados sobre \emph{todo} el conjunto de evaluación (y no promediados por imagen). Las métricas principales son el coeficiente Dice~\cite{dice1945measures} y el índice de Jaccard (IoU):
\begin{equation}\label{eq:dice}
  \DiceOp_c = \frac{2\,|P_c \cap G_c|}{|P_c| + |G_c|},
  \qquad
  \IoUOp_c = \frac{|P_c \cap G_c|}{|P_c \cup G_c|},
  \qquad
  \DiceOp_c = \frac{2\,\IoUOp_c}{1 + \IoUOp_c}.
\end{equation}
El \emph{Dice medio} es la media aritmética sobre las tres clases tumorales, excluido el fondo, cuya inclusión inflaría artificialmente cualquier métrica. Como métricas complementarias se emplean el Dice de \emph{tumor completo} (lesión frente a fondo, independientemente del tipo), la precisión, la sensibilidad (\emph{recall}), la razón entre falsos negativos y falsos positivos (FN/FP), el porcentaje de píxeles tumorales clasificados como fondo y el porcentaje de cortes en los que no se predice ningún píxel tumoral. En la comparación final, la incertidumbre se cuantifica mediante \emph{bootstrap} pareado por paciente~\cite{efron1993bootstrap}: 10\,000 remuestreos estratificados por tipo de tumor de los 46 pacientes de prueba.

\subsection{Protocolo experimental}\label{subsec:protocolo}

Todos los entrenamientos se ejecutaron en GPU (Kaggle) con PyTorch~\cite{paszke2019pytorch}, precisión mixta automática, lotes de 8 imágenes y semilla 42. En cada entrenamiento se conserva el punto de control con mayor Dice medio de validación. El conjunto de prueba se reservó para comparaciones planificadas de antemano y se utilizó tres veces a lo largo del proyecto: referencia manual frente a la configuración del AG, configuración del AG frente al modelo refinado y modelo refinado frente al modelo con transferencia. La \cref{tab:ejecuciones} enumera las ejecuciones citadas en el resto del documento; su configuración completa se recoge en el \cref{app:configuraciones}.

\begin{table}[htbp]
  \centering
  \caption{Ejecuciones de entrenamiento referenciadas en el documento.}
  \label{tab:ejecuciones}
  \begin{tabularx}{\linewidth}{llX}
    \toprule
    \textbf{Ejecución} & \textbf{Fase} & \textbf{Propósito} \\
    \midrule
    Referencia & 2--4 & Configuración manual de hiperparámetros (punto de partida). \\
    A          & 3--4 & Hiperparámetros encontrados por el algoritmo genético. \\
    B, C, D    & 5    & Diseño factorial pérdida $\times$ aumento sobre los hiperparámetros del AG. \\
    D-ext, S, I & 5   & Entrenamiento prolongado y aislamiento de los bloques espacial y de intensidad. \\
    E16, E     & 6    & Recocido coseno y pesos Focal Tversky recalibrados, con 16 y 32 filtros. \\
    F          & 7    & Complemento: codificador ResNet34 preentrenado en ImageNet. \\
    \bottomrule
  \end{tabularx}
\end{table}

\section{Nivel interno: la U-Net como problema de optimización}\label{sec:unet}

Esta sección estudia el entrenamiento de la red como lo que es en esencia: un algoritmo de optimización continua. Se formula el problema (\cref{subsec:erm}), se analiza cómo la arquitectura define el espacio de búsqueda y la geometría del paisaje (\cref{subsec:arquitectura}), se construyen las funciones objetivo sustitutas y se deriva su gradiente (\cref{subsec:perdidas}), se describe el cálculo exacto de ese gradiente por retropropagación (\cref{subsec:retropropagacion}) y, finalmente, se detallan el método de descenso y la programación de su tamaño de paso (\cref{subsec:adam,subsec:programacion}).

\subsection{Formulación: minimización del riesgo empírico}\label{subsec:erm}

Sea $\Dtrain = \{(\vx_n,\vy_n)\}_{n=1}^{N}$ el conjunto de entrenamiento, con $N=2004$ cortes de $\Omega = 256\times256$ píxeles. La red define una función $f_{\vtheta}:\R^{H\times W}\to\R^{C\times H\times W}$ que asigna a cada píxel $i$ un vector de $C=4$ \emph{logits} $\vz_i$, transformados en probabilidades mediante la función \emph{softmax}:
\begin{equation}\label{eq:softmax}
  p_{i,c} = \frac{\exp(z_{i,c})}{\sum_{k=0}^{C-1}\exp(z_{i,k})}, \qquad c \in \{0,\dots,C-1\}.
\end{equation}
El entrenamiento resuelve el problema de minimización del riesgo empírico regularizado
\begin{equation}\label{eq:erm}
  \min_{\vtheta\in\R^{d}} \; J(\vtheta) = \frac{1}{N}\sum_{n=1}^{N} \Loss\bigl(f_{\vtheta}(\vx_n),\vy_n\bigr) + \frac{\lambda}{2}\,\lVert\vtheta\rVert_2^2,
\end{equation}
donde $\Loss$ es la función de pérdida y $\lambda$ el coeficiente de decaimiento de pesos. Este problema tiene cuatro rasgos que lo distinguen de la optimización clásica y que justifican las técnicas empleadas:

\begin{enumerate}
  \item \textbf{Dimensión muy alta.} El vector de decisión $\vtheta$ tiene $d \approx 1.81\cdot10^{6}$ componentes para la red de 16 filtros y $d \approx 7.24\cdot10^{6}$ para la de 32. Cualquier método que requiera la matriz hessiana, de tamaño $d\times d$, es inviable; solo son practicables los métodos de primer orden.
  \item \textbf{No convexidad.} $J$ es la composición de decenas de transformaciones no lineales. Presenta múltiples mínimos locales, puntos de silla y simetrías (permutar los filtros de una capa no altera la función), de modo que no se busca el mínimo global sino un punto estacionario de buena calidad.
  \item \textbf{Estocasticidad.} Evaluar $J$ exactamente exige recorrer todo el conjunto. En su lugar, en cada iteración se estima el gradiente sobre un minilote $\mathcal{B}$ de $|\mathcal{B}|=8$ imágenes,
  \begin{equation}\label{eq:gradiente-estocastico}
    \hat{\vg}(\vtheta) = \frac{1}{|\mathcal{B}|}\sum_{n\in\mathcal{B}} \nabla_{\vtheta}\Loss\bigl(f_{\vtheta}(\vx_n),\vy_n\bigr) + \lambda\vtheta,
    \qquad \E_{\mathcal{B}}\bigl[\hat{\vg}(\vtheta)\bigr] = \nabla J(\vtheta),
  \end{equation}
  un estimador insesgado pero ruidoso. Con $N=2004$ y $|\mathcal{B}|=8$, cada época consta de 251 iteraciones de descenso.
  \item \textbf{La métrica de interés no es diferenciable.} El coeficiente Dice de la \cref{eq:dice} se calcula sobre la predicción discreta $\hat{y}_i=\argmax_c p_{i,c}$, que es constante a trozos en $\vtheta$: su gradiente es nulo en casi todo punto. Es imprescindible optimizar una \emph{función sustituta} diferenciable que se correlacione con el Dice (\cref{subsec:perdidas}).
\end{enumerate}

Además, el objetivo final no es minimizar $J$ sobre $\Dtrain$, sino generalizar a pacientes nuevos. Por ello, la trayectoria de optimización se supervisa con el conjunto de validación y se conserva el iterado con mayor Dice de validación, una forma de parada temprana que actúa como regularización implícita.

\subsection{La arquitectura como parametrización del espacio de búsqueda}\label{subsec:arquitectura}

La arquitectura determina qué funciones puede representar la red y, por tanto, el espacio en el que el optimizador busca. Se ha implementado una U-Net~\cite{ronneberger2015unet} con dos hiperparámetros estructurales, el número de filtros base $F$ y la profundidad $D$, que el algoritmo genético de la \cref{sec:ag} puede modificar sin alterar el código (\cref{fig:unet}).

\begin{itemize}
  \item \textbf{Codificador (contracción).} Consta de $D$ niveles. El nivel $\ell\in\{0,\dots,D-1\}$ aplica un bloque doble $(\text{Conv}\,3{\times}3 \to \text{BN} \to \text{ReLU})\times2$ con $2^{\ell}F$ canales y reduce después la resolución a la mitad mediante \emph{max pooling} $2\times2$. A medida que la resolución disminuye, el campo receptivo crece y los mapas de características codifican información más contextual.
  \item \textbf{Cuello de botella.} Un bloque doble con $2^{D}F$ canales a resolución $256/2^{D}$.
  \item \textbf{Decodificador (expansión).} Cada nivel duplica la resolución por interpolación bilineal, proyecta a $2^{\ell}F$ canales con una convolución $1\times1$, concatena las características del nivel homólogo del codificador y aplica un bloque doble.
  \item \textbf{Clasificador.} Una convolución $1\times1$ produce los $C=4$ logits por píxel.
\end{itemize}

\begin{figure}[htbp]
  \centering
  \begin{tikzpicture}[
      font=\small,
      enc/.style={draw=cazul!75!black, fill=cazul!20, rounded corners=1pt, inner sep=0pt},
      dec/.style={draw=cverde!70!black, fill=cverde!20, rounded corners=1pt, inner sep=0pt},
      bot/.style={draw=cpurpura!80!black, fill=cpurpura!25, rounded corners=1pt, inner sep=0pt},
      cls/.style={draw=cnaranja!80!black, fill=cnaranja!35, rounded corners=1pt, inner sep=0pt},
      lbl/.style={font=\scriptsize, align=center, text=black!80},
      pool/.style={-{Stealth[length=2.2mm]}, line width=0.9pt, draw=cbermellon},
      up/.style={-{Stealth[length=2.2mm]}, line width=0.9pt, draw=cverde!80!black},
      skip/.style={-{Stealth[length=2mm]}, line width=0.7pt, draw=cgris, dashed},
      flow/.style={-{Stealth[length=2mm]}, line width=0.7pt, draw=black!70},
      grad/.style={-{Stealth[length=2.4mm]}, line width=1pt, draw=cbermellon!85!black, densely dashed},
    ]
    \node[enc, minimum width=0.28cm, minimum height=1.50cm] (e0) at (1.0,  0.00) {};
    \node[enc, minimum width=0.40cm, minimum height=1.20cm] (e1) at (2.0, -1.45) {};
    \node[enc, minimum width=0.55cm, minimum height=0.95cm] (e2) at (3.0, -2.90) {};
    \node[enc, minimum width=0.72cm, minimum height=0.75cm] (e3) at (4.0, -4.35) {};
    \node[bot, minimum width=0.95cm, minimum height=0.60cm] (b)  at (5.4, -5.80) {};
    \node[dec, minimum width=0.72cm, minimum height=0.75cm] (d3) at (6.8, -4.35) {};
    \node[dec, minimum width=0.55cm, minimum height=0.95cm] (d2) at (7.8, -2.90) {};
    \node[dec, minimum width=0.40cm, minimum height=1.20cm] (d1) at (8.8, -1.45) {};
    \node[dec, minimum width=0.28cm, minimum height=1.50cm] (d0) at (9.8,  0.00) {};
    \node[cls, minimum width=0.20cm, minimum height=1.50cm] (c)  at (10.9, 0.00) {};

    \node[lbl, anchor=north east] at ($(e0.south west)+(-0.12,0)$) {$F$\\$256^2$};
    \node[lbl, anchor=north west] at ($(d0.south east)+(0.02,0)$) {$F$\\$256^2$};
    \node[lbl, left=3pt of e1] {$2F$\\$128^2$};
    \node[lbl, left=3pt of e2] {$4F$\\$64^2$};
    \node[lbl, left=3pt of e3] {$8F$\\$32^2$};
    \node[lbl, below=3pt of b] {$16F$, $16^2$\\cuello de botella};
    \node[lbl, right=3pt of d3] {$8F$\\$32^2$};
    \node[lbl, right=3pt of d2] {$4F$\\$64^2$};
    \node[lbl, right=3pt of d1] {$2F$\\$128^2$};
    \node[lbl, above=2pt of c] {Conv $1{\times}1$};

    \node[draw=black!60, fill=black!8, minimum size=0.9cm, align=center, font=\small] (x) at (-0.65,0) {$\vx$};
    \node[lbl, below=2pt of x] {$1\times256^2$};
    \node[align=center, right=0.55cm of c] (p) {$\vp$};
    \node[lbl, below=2pt of p] {$4\times256^2$\\softmax};
    \node[draw=cbermellon!85!black, fill=cbermellon!10, rounded corners=2pt, inner sep=4pt, right=0.55cm of p] (L) {$\Loss(\vp,\vy)$};
    \node[align=center, below=0.75cm of L, font=\small] (y) {$\vy$};
    \node[lbl, below=1pt of y] {máscara real};

    \draw[flow] (x.east) -- (e0.west);
    \draw[flow] (d0.east) -- (c.west);
    \draw[flow] (c.east) -- (p.west);
    \draw[flow] (p.east) -- (L.west);
    \draw[flow] (y.north) -- (L.south);

    \draw[pool] (e0.south) -- (e1.north);
    \draw[pool] (e1.south) -- (e2.north);
    \draw[pool] (e2.south) -- (e3.north);
    \draw[pool] (e3.south) -- (b.north west);
    \draw[up] (b.north east) -- (d3.south);
    \draw[up] (d3.north) -- (d2.south);
    \draw[up] (d2.north) -- (d1.south);
    \draw[up] (d1.north) -- (d0.south);

    \draw[skip] (e0.east) -- (d0.west);
    \draw[skip] (e1.east) -- (d1.west);
    \draw[skip] (e2.east) -- (d2.west);
    \draw[skip] (e3.east) -- (d3.west);

    \draw[grad] (L.north) -- ++(0,1.05) -| ($(e0.north)+(0,0.05)$);
    \node[font=\scriptsize, text=cbermellon!85!black, fill=white, inner sep=2pt] at (6.2,1.80) {retropropagación: $\nabla_{\vtheta}\Loss$ alcanza todas las capas a través de ambos caminos};

    \draw[decorate, decoration={brace, mirror, amplitude=5pt}, draw=black!60] (0.7,-5.25) -- (4.3,-5.25) node[midway, below=6pt, lbl] {codificador};
    \draw[decorate, decoration={brace, mirror, amplitude=5pt}, draw=black!60] (6.5,-5.25) -- (10.1,-5.25) node[midway, below=6pt, lbl] {decodificador};

    \node[enc, minimum width=0.25cm, minimum height=0.40cm] (le1) at (1.2,-7.3) {};
    \node[lbl, right=3pt of le1] (lt1) {bloque doble $(\text{Conv}\,3{\times}3\text{--BN--ReLU})\times2$};
    \coordinate[right=0.7cm of lt1] (la2);
    \coordinate[right=0.7cm of la2] (lb2);
    \draw[pool] (la2) -- (lb2);
    \node[lbl, right=3pt of lb2] {\emph{max pooling} $2{\times}2$};
    \coordinate (la3) at (1.05,-7.85);
    \coordinate[right=0.7cm of la3] (lb3);
    \draw[up] (la3) -- (lb3);
    \node[lbl, right=3pt of lb3] (lt3) {interpolación bilineal $\times2$ + Conv $1{\times}1$};
    \coordinate[right=0.7cm of lt3] (la4);
    \coordinate[right=0.7cm of la4] (lb4);
    \draw[skip] (la4) -- (lb4);
    \node[lbl, right=3pt of lb4] {conexión de salto (concatenación)};
  \end{tikzpicture}
  \caption{Arquitectura U-Net implementada, mostrada para la profundidad $D=4$ seleccionada por el algoritmo genético. Junto a cada bloque se indican el número de canales (en función de los filtros base $F\in\{8,16,32\}$) y la resolución espacial. La línea discontinua superior representa el flujo del gradiente durante la retropropagación: las conexiones de salto ofrecen a cada nivel del codificador un segundo camino, corto, por el que recibir la señal de error.}
  \label{fig:unet}
\end{figure}

Desde el punto de vista de la optimización, cada componente de la arquitectura cumple una función precisa:

\begin{itemize}
  \item \textbf{Conexiones de salto.} Si $\mathbf{h}_{\ell}$ es la salida del nivel $\ell$ del codificador, esta alimenta tanto al nivel inferior como, por concatenación, al nivel homólogo del decodificador. Por la regla de la cadena multivariable, su gradiente es la suma de ambas contribuciones,
  \begin{equation}\label{eq:gradiente-salto}
    \frac{\partial\Loss}{\partial\mathbf{h}_{\ell}} = \underbrace{\frac{\partial\Loss}{\partial\mathbf{h}_{\ell}}\bigg|_{\text{camino profundo}}}_{\text{atraviesa } 2(D-\ell) \text{ bloques}} + \underbrace{\frac{\partial\Loss}{\partial\mathbf{h}_{\ell}}\bigg|_{\text{salto}}}_{\text{camino corto}},
  \end{equation}
  de modo que los primeros niveles reciben una señal de error que no se atenúa al atravesar toda la red. Esto mitiga el desvanecimiento del gradiente y permite optimizar a la vez el contexto global y el detalle espacial fino que exige la delimitación de bordes.
  \item \textbf{Normalización por lotes.} Normaliza las preactivaciones con las estadísticas del minilote y las reescala con parámetros aprendidos~\cite{ioffe2015batchnorm}. Esta reparametrización mejora el condicionamiento del problema y tolera tasas de aprendizaje mayores. Las convoluciones prescinden del sesgo, que queda absorbido por el desplazamiento de la normalización.
  \item \textbf{Inicialización de He.} Los pesos convolucionales se inicializan con varianza $2/n_{\mathrm{in}}$~\cite{he2015delving}, lo que conserva la varianza de las activaciones a través de las capas ReLU. El punto de partida del descenso se sitúa así en una región donde los gradientes ni explotan ni se anulan.
  \item \textbf{Interpolación bilineal con proyección $1\times1$.} Sustituye a la convolución traspuesta y evita los artefactos en damero que esta puede dejar en la máscara.
\end{itemize}

El número de parámetros crece aproximadamente como $F^{2}$ y como $4^{D}$, porque cada convolución $3\times3$ entre $a$ y $b$ canales aporta $9ab$ pesos. Para $D=4$, la red tiene $1\,811\,780$ parámetros con $F=16$ y $7\,240\,324$ con $F=32$, exactamente cuatro veces más. La \cref{tab:parametros} de la \cref{sec:ag} muestra el tamaño de todas las arquitecturas del espacio de búsqueda.

\subsection{Funciones de pérdida: objetivos sustitutos diferenciables}\label{subsec:perdidas}

Sean $\vp_{c}$ y $\vg_{c}$ los mapas de probabilidad y la codificación \emph{one-hot} de la clase $c$ sobre todos los píxeles del minilote. Todas las pérdidas de región se calculan con agregación a nivel de minilote, excluyen el fondo y promedian solo sobre el conjunto $\mathcal{C}_{\mathcal{B}}$ de clases tumorales presentes en el minilote. Este último detalle es necesario porque cada corte contiene un único tipo de tumor: una clase ausente del minilote no aporta información de solapamiento y, si se incluyera, contaría como una predicción perfecta que diluiría la pérdida.

\paragraph{Dice suave.} Se obtiene sustituyendo las predicciones discretas por probabilidades en la \cref{eq:dice}~\cite{milletari2016vnet}:
\begin{equation}\label{eq:dice-suave}
  \Loss_{\mathrm{Dice}} = 1 - \frac{1}{|\mathcal{C}_{\mathcal{B}}|}\sum_{c\in\mathcal{C}_{\mathcal{B}}} \frac{2\sum_i p_{i,c}\,g_{i,c} + \varepsilon}{\sum_i p_{i,c} + \sum_i g_{i,c} + \varepsilon}, \qquad \varepsilon = 10^{-5}.
\end{equation}
Al ser un cociente de sumas sobre todos los píxeles, es insensible a la proporción entre fondo y tumor: mide solapamiento relativo y no cuenta los verdaderos negativos, que son abrumadoramente mayoritarios.

\paragraph{Pérdida focal.} Es una entropía cruzada reponderada~\cite{lin2017focal}:
\begin{equation}\label{eq:focal}
  \Loss_{\mathrm{Focal}} = -\frac{1}{|\Omega_{\mathcal{B}}|}\sum_{i} \bigl(1-p_{i,y_i}\bigr)^{\gamma}\,\log p_{i,y_i}, \qquad \gamma = 2.
\end{equation}
El factor modulador $(1-p_{i,y_i})^{\gamma}$ reduce la contribución de los píxeles ya bien clasificados: un píxel de fondo con $p=0.9$ pesa $0.01$, cien veces menos que en la entropía cruzada (el caso $\gamma=0$). Como la inmensa mayoría de los píxeles son fondo fácil, el gradiente se concentra en los píxeles difíciles, típicamente los bordes de la lesión.

\paragraph{Pérdida Focal Tversky.} El índice de Tversky~\cite{tversky1977features,salehi2017tversky} generaliza el Dice ponderando por separado los dos tipos de error. Con recuentos suaves $\mathrm{TP}_c=\sum_i p_{i,c}g_{i,c}$, $\mathrm{FN}_c=\sum_i(1-p_{i,c})g_{i,c}$ y $\mathrm{FP}_c=\sum_i p_{i,c}(1-g_{i,c})$,
\begin{equation}\label{eq:tversky}
  \mathrm{TI}_c = \frac{\mathrm{TP}_c}{\mathrm{TP}_c + \alpha\,\mathrm{FN}_c + \beta\,\mathrm{FP}_c},
  \qquad
  \Loss_{\mathrm{FT}} = \frac{1}{|\mathcal{C}_{\mathcal{B}}|}\sum_{c\in\mathcal{C}_{\mathcal{B}}} \bigl(1-\mathrm{TI}_c\bigr)^{\gamma_{T}}.
\end{equation}
Con $\alpha=\beta=0.5$ y $\gamma_T=1$ se recupera el Dice suave. La variante focal~\cite{abraham2019focaltversky} con $\gamma_T=4/3$ concentra el esfuerzo en las clases peor segmentadas.

\paragraph{Combinación por penalización.} Las pérdidas de región optimizan directamente el solapamiento, pero su gradiente es inestable al principio del entrenamiento, cuando el solapamiento es casi nulo. La pérdida focal aporta un gradiente suave y local en cada píxel. Ambas se combinan por escalarización ponderada, el método clásico de penalización para objetivos múltiples:
\begin{equation}\label{eq:combinada}
  \Loss = w_{r}\,\Loss_{\text{región}} + w_{p}\,\Loss_{\mathrm{Focal}}, \qquad w_r = w_p = 0.5,
\end{equation}
donde $\Loss_{\text{región}}$ es $\Loss_{\mathrm{Dice}}$ (pérdida \emph{Dice-Focal}, Fases 2--4) o $\Loss_{\mathrm{FT}}$ (pérdida \emph{Focal Tversky-Focal}, Fases 5--7).

\subsubsection{Análisis del gradiente de la pérdida Focal Tversky}\label{subsubsec:gradiente-tversky}

El comportamiento de la pérdida durante el descenso se entiende derivándola. Sea $G_c=\sum_i g_{i,c}$ el número de píxeles reales de la clase, constante respecto a $\vtheta$, de modo que $\mathrm{FN}_c = G_c - \mathrm{TP}_c$, y sea $S_c = \mathrm{TP}_c + \alpha\mathrm{FN}_c + \beta\mathrm{FP}_c$ el denominador de la \cref{eq:tversky}. Derivando respecto a la probabilidad de un píxel se obtienen dos casos:
\begin{equation}\label{eq:gradiente-ti}
  \frac{\partial\,\mathrm{TI}_c}{\partial p_{i,c}} =
  \begin{cases}
    \dfrac{\alpha\,G_c + \beta\,\mathrm{FP}_c}{S_c^{2}} > 0 & \text{si } g_{i,c}=1 \quad\text{(píxel tumoral)},\\[3mm]
    -\dfrac{\beta\,\mathrm{TP}_c}{S_c^{2}} < 0 & \text{si } g_{i,c}=0 \quad\text{(píxel de fondo)}.
  \end{cases}
\end{equation}
La lectura es directa. El incentivo para elevar la probabilidad de un píxel tumoral no detectado crece con $\alpha$, y la penalización por marcar un píxel de fondo crece con $\beta$. Elegir $\alpha>\beta$ inclina, por tanto, el óptimo hacia una mayor sensibilidad a costa de precisión. Esta relación analítica es la base cuantitativa de las decisiones de las Fases 5 y 6 (\cref{sec:refinamiento}), en las que $\alpha$ y $\beta$ se calibran a partir de la razón FN/FP observada.

El exponente focal actúa a nivel de clase. Por la regla de la cadena,
\begin{equation}\label{eq:gradiente-focal-tversky}
  \frac{\partial\Loss_{\mathrm{FT}}}{\partial\,\mathrm{TI}_c} = -\frac{\gamma_T}{|\mathcal{C}_{\mathcal{B}}|}\,\bigl(1-\mathrm{TI}_c\bigr)^{\gamma_T-1}.
\end{equation}
Con $\gamma_T=4/3>1$, el factor $(1-\mathrm{TI}_c)^{1/3}$ crece con el error. Una clase con $\mathrm{TI}_c\approx0.6$, como el glioma, recibe $(0.40/0.11)^{1/3}\approx1.5$ veces el gradiente de una clase con $\mathrm{TI}_c\approx0.89$, como el meningioma. La pérdida redirige automáticamente el esfuerzo de optimización hacia la clase cuello de botella.

Por último, el gradiente se propaga a los logits a través del jacobiano de la \emph{softmax}, $\partial p_{i,c}/\partial z_{i,k} = p_{i,c}(\delta_{ck}-p_{i,k})$, de modo que
\begin{equation}\label{eq:gradiente-logits}
  \frac{\partial\Loss}{\partial z_{i,k}} = \sum_{c} \frac{\partial\Loss}{\partial p_{i,c}}\; p_{i,c}\bigl(\delta_{ck}-p_{i,k}\bigr),
\end{equation}
que es el vector de error inicial a partir del cual comienza la retropropagación.

\subsection{Retropropagación: cálculo exacto del gradiente}\label{subsec:retropropagacion}

La red es un grafo acíclico dirigido de operaciones diferenciables (o diferenciables a trozos). Escribiendo cada capa como $\mathbf{a}^{(l)} = f_l(\mathbf{a}^{(l-1)};\vtheta_l)$ y definiendo el error $\bm{\delta}^{(l)} = \partial\Loss/\partial\mathbf{a}^{(l)}$, la retropropagación~\cite{rumelhart1986backprop} aplica recursivamente la regla de la cadena en sentido inverso:
\begin{equation}\label{eq:backprop}
  \bm{\delta}^{(l-1)} = \Bigl(\frac{\partial f_l}{\partial\mathbf{a}^{(l-1)}}\Bigr)^{\!\top}\bm{\delta}^{(l)},
  \qquad
  \nabla_{\vtheta_l}\Loss = \Bigl(\frac{\partial f_l}{\partial\vtheta_l}\Bigr)^{\!\top}\bm{\delta}^{(l)}.
\end{equation}
Nunca se forman jacobianos explícitos: cada operación implementa su producto jacobiano-vector traspuesto. En una convolución, la propagación del error equivale a una convolución con el núcleo invertido, y el gradiente del núcleo es la correlación entre las activaciones de entrada y el error de salida. En la ReLU, el error se multiplica por el indicador $\mathbb{1}[\mathbf{a}>0]$. En el \emph{max pooling}, el error se encamina hacia la posición que contenía el máximo (un subgradiente). Cuando un tensor alimenta a varios consumidores, como en las conexiones de salto, los errores se suman (\cref{eq:gradiente-salto}).

La propiedad decisiva para la optimización es su coste: una pasada hacia atrás cuesta del orden de dos pasadas hacia delante, \emph{independientemente de} $d$~\cite{goodfellow2016deep}. Estimar el mismo gradiente por diferencias finitas exigiría $d+1\approx7\cdot10^{6}$ evaluaciones de la red. La retropropagación es lo que hace factible el descenso de gradiente en millones de variables.

Para acelerar el cálculo se emplea precisión mixta~\cite{micikevicius2018mixed}: la pasada hacia delante se ejecuta en media precisión donde es seguro, pero la pérdida se evalúa en precisión simple, ya que la constante $\varepsilon=10^{-5}$ de las \cref{eq:dice-suave,eq:tversky} es subnormal en media precisión y perdería exactitud. Para evitar que los gradientes pequeños se anulen, se retropropaga $s\cdot\Loss$ con un factor de escala dinámico $s$ y los gradientes se dividen por $s$ antes de actualizar los pesos.

\subsection{Método de descenso: Adam con regularización \texorpdfstring{$L_2$}{L2}}\label{subsec:adam}

El descenso de gradiente estocástico~\cite{robbins1951stochastic} actualiza $\vtheta_{t+1}=\vtheta_t-\eta\,\hat{\vg}_t$. En este trabajo se usa su variante adaptativa Adam~\cite{kingma2015adam}, que en la iteración $t$ calcula
\begin{equation}\label{eq:adam}
  \begin{aligned}
    \vg_t &= \hat{\vg}(\vtheta_{t-1}) \quad\text{(incluye } \lambda\vtheta_{t-1}\text{)}, &
    \mathbf{m}_t &= \beta_1\mathbf{m}_{t-1} + (1-\beta_1)\,\vg_t, &
    \mathbf{v}_t &= \beta_2\mathbf{v}_{t-1} + (1-\beta_2)\,\vg_t^{2},\\
    \hat{\mathbf{m}}_t &= \frac{\mathbf{m}_t}{1-\beta_1^{t}}, &
    \hat{\mathbf{v}}_t &= \frac{\mathbf{v}_t}{1-\beta_2^{t}}, &
    \vtheta_t &= \vtheta_{t-1} - \eta\,\frac{\hat{\mathbf{m}}_t}{\sqrt{\hat{\mathbf{v}}_t}+\epsilon},
  \end{aligned}
\end{equation}
con $\beta_1=0.9$, $\beta_2=0.999$ y $\epsilon=10^{-8}$, y con todas las operaciones vectoriales componente a componente. El primer momento $\mathbf{m}_t$ es un término de inercia que promedia el ruido del gradiente estocástico. El segundo momento $\mathbf{v}_t$ actúa como un precondicionador diagonal que normaliza la escala de cada coordenada. El resultado es que el desplazamiento de cada peso por iteración queda acotado aproximadamente por $\eta$, lo que convierte a la \emph{tasa de aprendizaje} en el hiperparámetro más influyente del proceso: fija el tamaño del paso en el espacio de parámetros.

El decaimiento de pesos se implementa, como en PyTorch, sumando $\lambda\vtheta$ al gradiente; es decir, como el gradiente de la penalización $L_2$ de la \cref{eq:erm}. Como este término atraviesa después la normalización adaptativa, su efecto regularizador es más débil y menos uniforme que el del decaimiento desacoplado~\cite{loshchilov2019adamw}. Este matiz ayuda a interpretar la elección del algoritmo genético en la \cref{subsec:resultados-ag}.

En suma, el algoritmo de optimización interno queda completamente determinado por cuatro magnitudes que no se aprenden por gradiente: la tasa de aprendizaje $\eta$ y el decaimiento $\lambda$, que gobiernan el método de descenso, y los filtros $F$ y la profundidad $D$, que definen el espacio de búsqueda. Estas cuatro magnitudes forman el cromosoma del algoritmo genético de la \cref{sec:ag}.

\subsection{Programación de la tasa de aprendizaje}\label{subsec:programacion}

Con una tasa constante, el descenso estocástico no converge a un punto: los iterados oscilan en un entorno del punto estacionario cuyo radio crece con $\eta$ y con la varianza del gradiente. En la práctica, esto se traduce en fluctuaciones del Dice de validación de una época a la siguiente, que convierten la selección del mejor punto de control en un sorteo. El \emph{recocido coseno}~\cite{loshchilov2017sgdr}, en su versión sin reinicios, reduce la tasa a lo largo del entrenamiento según
\begin{equation}\label{eq:coseno}
  \eta_e = \frac{\eta_0}{2}\Bigl(1+\cos\frac{\pi e}{T}\Bigr), \qquad e = 0,1,\dots,T-1,
\end{equation}
donde $e$ es la época y $T$ el número total de épocas. Se dan pasos grandes al principio, que exploran distintas cuencas del paisaje, y pasos cada vez menores al final, que reducen el ruido y permiten asentarse en el fondo de la cuenca. La analogía con la temperatura del recocido simulado es deliberada: en ambos casos, un parámetro que decrece con el tiempo regula la transición de la exploración a la explotación. Esta técnica se incorpora en la Fase 6 (\cref{sec:refinamiento}).

\subsection{Comportamiento observado del proceso de optimización}\label{subsec:comportamiento}

La \cref{fig:perdidas} muestra la evolución de la función objetivo para dos ejecuciones que comparten arquitectura ($F=16$, $D=4$), pérdida (Dice-Focal), optimizador y datos, y que solo difieren en los hiperparámetros del optimizador: la configuración manual de referencia ($\eta=10^{-3}$, $\lambda=10^{-4}$) y la encontrada por el algoritmo genético ($\eta=3.15\cdot10^{-4}$, $\lambda=1.91\cdot10^{-6}$).

\begin{figure}[htbp]
  \centering
  \begin{tikzpicture}
    \begin{axis}[
        informe,
        height=6.4cm,
        xlabel={Época},
        ylabel={Pérdida Dice-Focal},
        xmin=0, xmax=100, ymin=0, ymax=0.6,
        legend pos=north east,
        legend columns=2,
      ]
      \addplot[cazul, line width=1.1pt] table[x=epoca, y=perdida_train, col sep=comma] {data/historial_base.csv};
      \addlegendentry{Referencia, entrenamiento}
      \addplot[cazul, line width=1.1pt, densely dashed] table[x=epoca, y=perdida_val, col sep=comma] {data/historial_base.csv};
      \addlegendentry{Referencia, validación}
      \addplot[cbermellon, line width=1.1pt] table[x=epoca, y=perdida_train, col sep=comma] {data/historial_A.csv};
      \addlegendentry{AG (A), entrenamiento}
      \addplot[cbermellon, line width=1.1pt, densely dashed] table[x=epoca, y=perdida_val, col sep=comma] {data/historial_A.csv};
      \addlegendentry{AG (A), validación}
    \end{axis}
  \end{tikzpicture}
  \caption{Evolución de la función objetivo para la configuración manual de referencia (40 épocas) y la configuración del algoritmo genético (100 épocas), con idéntica arquitectura y pérdida. La configuración del AG, con una tasa de aprendizaje tres veces menor, desciende más rápido: en la época 10, su pérdida de entrenamiento es 0.188 frente a 0.282.}
  \label{fig:perdidas}
\end{figure}

La figura ilustra tres propiedades del proceso. Primero, el descenso es eficaz: la pérdida de entrenamiento cae desde $\approx0.55$ hasta $0.18$ (referencia, época 40) y $0.08$ (configuración del AG, época 100). Segundo, es estocástico: la pérdida de entrenamiento aumenta respecto a la época anterior en 10 de las 39 transiciones de la referencia y en 40 de las 99 de la configuración del AG, y la de validación oscila aún más. Tercero, y más importante para este trabajo, la relación entre hiperparámetros y convergencia no es intuitiva. Una tasa de aprendizaje \emph{menor}, combinada con un decaimiento de pesos casi nulo, produce un descenso \emph{más rápido}: en la época 40, la configuración del AG alcanza una pérdida de entrenamiento de 0.119 frente a 0.180 y una pérdida mínima de validación de 0.184 frente a 0.228. La interpretación más plausible es que el paso de $10^{-3}$ es excesivo para este paisaje y provoca oscilaciones que frenan el progreso, aunque parte de la diferencia puede deberse a la penalización $L_2$ más intensa de la referencia; con una sola ejecución de cada configuración no es posible separar ambos efectos.

Si con la misma red, la misma función objetivo y el mismo optimizador dos elecciones de $(\eta,\lambda)$ conducen a trayectorias y puntos finales tan distintos, y si esa relación no es monótona ni predecible a priori, la elección de los hiperparámetros es en sí misma un problema de optimización. A ese segundo nivel se dedica la siguiente sección.

\section{Nivel externo: metaoptimización de hiperparámetros con un algoritmo genético}\label{sec:ag}

La \cref{sec:unet} ha mostrado que el algoritmo de optimización interno queda determinado por cuatro magnitudes que el descenso de gradiente no puede ajustar, y que la calidad de la solución depende de ellas de forma no monótona. Esta sección aborda el problema externo de la \cref{eq:binivel} con un algoritmo genético implementado desde cero.

\subsection{Caracterización del problema de optimización de hiperparámetros}\label{subsec:caracterizacion}

La función que se desea maximizar, $\vlambda \mapsto \DiceOp_{\mathrm{val}}(\vtheta^{\star}(\vlambda))$, reúne todas las propiedades que hacen inviables los métodos clásicos~\cite{feurer2019hpo}:

\begin{enumerate}
  \item \textbf{Caja negra.} No tiene expresión analítica ni gradiente respecto a $\vlambda$: cada evaluación es la salida de un proceso de entrenamiento completo.
  \item \textbf{Costosa.} Un entrenamiento de 100 épocas dura en torno a 98 minutos en GPU, lo que limita el presupuesto a unas pocas decenas de evaluaciones.
  \item \textbf{Ruidosa.} El resultado depende de la inicialización aleatoria y del orden de los minilotes.
  \item \textbf{Mixta y con escalas heterogéneas.} Combina variables continuas que abarcan varios órdenes de magnitud ($\eta$, $\lambda$) con variables discretas ordinales ($F$, $D$).
  \item \textbf{Multimodal y con interacciones.} El valor óptimo de un hiperparámetro depende de los demás; por ejemplo, la tasa de aprendizaje adecuada cambia con el tamaño de la red.
\end{enumerate}

La búsqueda en rejilla crece exponencialmente con la dimensión y malgasta evaluaciones en las dimensiones poco influyentes. La búsqueda aleatoria es una referencia sólida~\cite{bergstra2012random}, pero no aprovecha la información de las evaluaciones previas. Las metaheurísticas poblacionales, como los algoritmos genéticos~\cite{holland1975adaptation,goldberg1989genetic} o la optimización por enjambre de partículas~\cite{kennedy1995pso}, sí lo hacen: concentran progresivamente la búsqueda en las regiones prometedoras sin requerir gradientes. Se eligió un algoritmo genético por tres razones: admite con naturalidad variables mixtas mediante una codificación adecuada, la selección basada en el orden es robusta frente al ruido de la aptitud, y su aplicación a hiperparámetros de redes profundas está bien establecida~\cite{young2015optimizing}.

\subsubsection{Por qué no se adoptó la búsqueda aleatoria}\label{subsubsec:aleatoria}

La búsqueda aleatoria es la referencia natural de cualquier método de optimización de hiperparámetros. Bergstra y Bengio~\cite{bergstra2012random} demostraron que, a igual presupuesto, supera a la búsqueda en rejilla, porque no malgasta evaluaciones en las dimensiones poco influyentes. En este trabajo se descartó como método de búsqueda por una razón de coste. Cada evaluación exige entrenar una red durante tres épocas, unos tres minutos de GPU, lo que limitó el presupuesto total a 46 evaluaciones. En ese régimen, el criterio decisivo es la eficiencia por muestra, y la búsqueda aleatoria tiene una limitación estructural: es una búsqueda no informada y sin memoria. Cada configuración se muestrea con independencia de las anteriores, de modo que lo aprendido en las evaluaciones ya realizadas no influye en dónde se evalúa a continuación.

Esta limitación puede cuantificarse. La probabilidad de que al menos una de $n$ muestras independientes y uniformes caiga en una región que ocupa una fracción $q$ del espacio de búsqueda es
\begin{equation}\label{eq:aleatoria}
  P_n(q) = 1-(1-q)^{n}.
\end{equation}
Con $n=46$, la búsqueda aleatoria alcanza casi con seguridad el 5\,\% mejor del espacio ($P_{46}(0.05)=0.91$), pero solo alcanza el 1\,\% mejor en el 37\,\% de los casos y el 0.2\,\% mejor en el 9\,\%. Localizar una región buena es, por tanto, sencillo; refinarla dentro del presupuesto no lo es, porque la búsqueda aleatoria sigue repartiendo las evaluaciones de manera uniforme incluso después de haber identificado las zonas prometedoras.

El algoritmo genético está diseñado precisamente para superar esta limitación. Su población actúa como una memoria de las evaluaciones pasadas, y sus operadores regulan de forma explícita el equilibrio entre dos fuerzas. La selección por torneo y el cruce \emph{explotan} esa memoria: concentran los nuevos muestreos alrededor de los individuos más aptos y recombinan sus componentes. La mutación gaussiana, junto con la componente extrapolante del cruce BLX-$\alpha$, mantiene la \emph{exploración} y evita la convergencia prematura. Con un presupuesto tan restringido, este equilibrio permite dedicar la mayor parte de las evaluaciones a refinar las regiones prometedoras sin abandonar el resto del espacio. La ventaja presupone que el paisaje de la aptitud tiene una estructura explotable, es decir, que configuraciones próximas obtienen aptitudes parecidas; el paisaje observado (\cref{fig:ag-paisaje}) muestra que así ocurre, y la \cref{subsec:resultados-ag} cuantifica cómo el AG concentró su presupuesto. Se trata, no obstante, de un argumento teórico: establece por qué cabe esperar que el AG sea más eficiente por muestra, no que lo haya sido en este problema concreto, cuestión que se discute en la \cref{subsec:validacion-ag}.

\subsection{Espacio de búsqueda y codificación del cromosoma}\label{subsec:codificacion}

Cada individuo se representa mediante un \emph{genotipo} real normalizado $\vu=(u_1,u_2,u_3,u_4)\in[0,1]^4$, que una función de decodificación $\varphi$ transforma en el \emph{fenotipo} $\vlambda=\varphi(\vu)=(\eta,\lambda,F,D)$. La \cref{tab:genes} detalla cada gen.

\begin{table}[htbp]
  \centering
  \caption{Codificación del cromosoma. La decodificación de los genes continuos es logarítmica y la de los categóricos asigna intervalos de igual anchura a cada valor. Las dos últimas columnas muestran el fenotipo y, entre paréntesis, el gen de la configuración manual de referencia y de la mejor solución encontrada.}
  \label{tab:genes}
  \small
  \setlength{\tabcolsep}{4pt}
  \begin{tabular}{clllll}
    \toprule
    \textbf{Gen} & \textbf{Hiperparámetro} & \textbf{Dominio} & \textbf{Decodificación} & \textbf{Referencia} & \textbf{Mejor (AG)} \\
    \midrule
    $u_1$ & Tasa de aprendizaje $\eta$ & $[10^{-4},10^{-2}]$ & $10^{-4+2u_1}$ & $10^{-3}$ (0.500) & $3.15\cdot10^{-4}$ (0.249) \\
    $u_2$ & Decaimiento de pesos $\lambda$ & $[10^{-6},10^{-3}]$ & $10^{-6+3u_2}$ & $10^{-4}$ (0.667) & $1.91\cdot10^{-6}$ (0.094) \\
    $u_3$ & Filtros base $F$ & $\{8,16,32\}$ & $(8,16,32)_{\lfloor 3u_3\rfloor}$ & 16 (0.500) & 16 (0.478) \\
    $u_4$ & Profundidad $D$ & $\{3,4,5\}$ & $(3,4,5)_{\lfloor 3u_4\rfloor}$ & 4 (0.500) & 4 (0.520) \\
    \bottomrule
  \end{tabular}
\end{table}

Las decisiones de diseño de esta codificación son las siguientes:

\begin{itemize}
  \item \textbf{Escala logarítmica para los genes continuos.} La tasa de aprendizaje y el decaimiento actúan de forma multiplicativa: pasar de $10^{-4}$ a $2\cdot10^{-4}$ tiene un efecto comparable a pasar de $10^{-3}$ a $2\cdot10^{-3}$. Con la decodificación $10^{a+(b-a)u}$, un gen uniforme en $[0,1]$ explora por igual cada orden de magnitud y los operadores trabajan sobre el exponente, que es la escala relevante.
  \item \textbf{Intervalos ordenados para los genes categóricos.} El intervalo $[0,1]$ se divide en tres tramos de anchura $1/3$, asociados a los valores en orden creciente (el índice se acota con $\min(\lfloor 3u\rfloor,2)$ para incluir $u=1$). Como las categorías son ordinales, los operadores aritméticos conservan su sentido: el cruce de un individuo de 8 filtros con otro de 32 puede producir uno de 16. Un valor conocido se codifica en el centro de su tramo, $(k+0.5)/3$, como se hizo con la configuración de referencia.
  \item \textbf{Redundancia controlada.} La correspondencia genotipo-fenotipo es sobreyectiva y no inyectiva en los genes categóricos: pequeñas variaciones de $u_3$ o $u_4$ dentro de un tramo son mutaciones neutras que no cambian la red. Esto permite la deriva genética sin coste de evaluación, gracias a la memoria descrita en la \cref{subsec:aptitud}.
\end{itemize}

El espacio de búsqueda abarca nueve arquitecturas cuyo tamaño varía en más de dos órdenes de magnitud (\cref{tab:parametros}), combinadas con un continuo bidimensional de configuraciones del optimizador. Todas las profundidades son compatibles con la entrada de $256\times256$ píxeles, que es divisible por $2^5$. La \cref{fig:cromosoma} ilustra la decodificación del mejor individuo encontrado.

\begin{table}[htbp]
  \centering
  \caption{Número de parámetros entrenables, en millones, de cada arquitectura del espacio de búsqueda.}
  \label{tab:parametros}
  \begin{tabular}{lccc}
    \toprule
    & $D=3$ & $D=4$ & $D=5$ \\
    \midrule
    $F=8$  & 0.11 & 0.45 & 1.82 \\
    $F=16$ & 0.45 & \textbf{1.81} & 7.25 \\
    $F=32$ & 1.80 & 7.24 & 29.00 \\
    \bottomrule
  \end{tabular}
\end{table}

\begin{figure}[htbp]
  \centering
  \begin{tikzpicture}[
      font=\small,
      gen/.style={draw=cazul!70!black, fill=fondogen, rounded corners=2pt, minimum width=2.7cm, minimum height=0.8cm},
      cat/.style={draw=cnaranja!80!black, fill=fondocat, rounded corners=2pt, minimum width=2.7cm, minimum height=0.8cm},
      head/.style={font=\scriptsize, text=black!75},
      fila/.style={font=\scriptsize\bfseries, text=black!70, align=flush left, anchor=west, text width=2.25cm},
      dec/.style={-{Stealth[length=2mm]}, line width=0.6pt, draw=black!60},
      marca/.style={draw=cbermellon, line width=1.4pt},
      res/.style={font=\small, align=flush center},
    ]
    \node[fila] at (-0.4, 0.0) {Genotipo\\$\vu\in[0,1]^4$};
    \node[fila] at (-0.4,-1.75) {Decodificación\\$\varphi(\vu)$};
    \node[fila] at (-0.4,-4.0) {Evaluación\\de la aptitud};

    \node[gen] (g1) at (3.4,0) {$u_1 = 0.249$};
    \node[gen] (g2) at (6.5,0) {$u_2 = 0.094$};
    \node[cat] (g3) at (9.6,0) {$u_3 = 0.478$};
    \node[cat] (g4) at (12.7,0) {$u_4 = 0.520$};
    \node[head, above=1pt of g1] {gen continuo: $\eta$};
    \node[head, above=1pt of g2] {gen continuo: $\lambda$};
    \node[head, above=1pt of g3] {gen categórico: $F$};
    \node[head, above=1pt of g4] {gen categórico: $D$};

    \foreach \cx/\u/\lo/\hi in {3.4/0.249/{10^{-4}}/{10^{-2}}, 6.5/0.094/{10^{-6}}/{10^{-3}}} {
      \draw[line width=0.8pt, draw=cazul!70!black] (\cx-1.2,-1.6) -- (\cx+1.2,-1.6);
      \draw[draw=cazul!70!black] (\cx-1.2,-1.5) -- (\cx-1.2,-1.7);
      \draw[draw=cazul!70!black] (\cx+1.2,-1.5) -- (\cx+1.2,-1.7);
      \node[font=\tiny, below=1pt] at (\cx-1.2,-1.7) {$\lo$};
      \node[font=\tiny, below=1pt] at (\cx+1.2,-1.7) {$\hi$};
      \draw[marca] (\cx-1.2+2.4*\u,-1.42) -- (\cx-1.2+2.4*\u,-1.78);
    }
    \foreach \cx/\u/\a/\b/\c in {9.6/0.478/8/16/32, 12.7/0.520/3/4/5} {
      \draw[draw=cnaranja!80!black, fill=fondocat] (\cx-1.2,-1.78) rectangle (\cx+1.2,-1.42);
      \draw[draw=cnaranja!80!black] (\cx-0.4,-1.78) -- (\cx-0.4,-1.42);
      \draw[draw=cnaranja!80!black] (\cx+0.4,-1.78) -- (\cx+0.4,-1.42);
      \node[font=\scriptsize] at (\cx-0.8,-1.6) {\a};
      \node[font=\scriptsize] at (\cx,-1.6) {\b};
      \node[font=\scriptsize] at (\cx+0.8,-1.6) {\c};
      \draw[marca] (\cx-1.2+2.4*\u,-1.32) -- (\cx-1.2+2.4*\u,-1.88);
    }

    \node[res] (f1) at (3.4,-2.6) {$\eta = 10^{-4+2u_1}$\\$= 3.15\cdot10^{-4}$};
    \node[res] (f2) at (6.5,-2.6) {$\lambda = 10^{-6+3u_2}$\\$= 1.91\cdot10^{-6}$};
    \node[res] (f3) at (9.6,-2.6) {$\lfloor 3u_3\rfloor = 1$\\$F = 16$};
    \node[res] (f4) at (12.7,-2.6) {$\lfloor 3u_4\rfloor = 1$\\$D = 4$};

    \foreach \i in {1,2,3,4} {
      \draw[dec] (g\i.south) -- ++(0,-0.48);
    }

    \node[draw=cverde!70!black, fill=cverde!10, rounded corners=3pt, text width=11.4cm, minimum height=1.0cm, align=flush center, inner sep=5pt] (ev) at (8.05,-4.0)
      {construir U-Net$(F,D)$ y Adam$(\eta,\lambda)$ $\;\longrightarrow\;$ entrenar 3 épocas con semilla fija\\[2pt] $\longrightarrow\;\; \aptitud(\vu) = \max_{t\le3} \DiceOp_{\mathrm{val}}^{(t)} = 0.5085$};
    \foreach \i in {1,2,3,4} {
      \draw[dec] (f\i.south) -- (f\i.south |- ev.north);
    }
  \end{tikzpicture}
  \caption{Codificación y decodificación del cromosoma, ilustradas con el mejor individuo encontrado. Los genes continuos (azul) se proyectan sobre una escala logarítmica y los categóricos (naranja) sobre tres tramos ordenados de igual anchura; la marca roja indica la posición del gen. El fenotipo resultante instancia la red y el optimizador cuya calidad mide la función de aptitud.}
  \label{fig:cromosoma}
\end{figure}

\subsection{Función de aptitud}\label{subsec:aptitud}

La aptitud de un genotipo es una aproximación de baja fidelidad del objetivo externo de la \cref{eq:binivel}:
\begin{equation}\label{eq:aptitud}
  \aptitud(\vu) =
  \begin{cases}
    \displaystyle\max_{t\in\{1,2,3\}} \DiceOp_{\mathrm{val}}\bigl(\vtheta_t(\varphi(\vu))\bigr) & \text{si el entrenamiento concluye},\\[2mm]
    0 & \text{si diverge (pérdida no finita) o agota la memoria},
  \end{cases}
\end{equation}
donde $\vtheta_t(\vlambda)$ son los pesos tras $t$ épocas de Adam con los hiperparámetros $\vlambda$, la pérdida Dice-Focal ($w_r=w_p=0.5$, $\gamma=2$), el aumento básico, lotes de 8 imágenes y precisión mixta. Sus propiedades son las siguientes:

\begin{itemize}
  \item \textbf{Baja fidelidad.} Tres épocas equivalen a 753 iteraciones de descenso. Cada evaluación duró de media 178.5\,s (entre 146 y 223\,s), frente a unos 98 minutos de un entrenamiento de 100 épocas: un coste unas 33 veces menor. El coste crece débilmente con el tamaño de la red (161\,s con $F=8$, 179\,s con $F=16$ y 198\,s con $F=32$), porque el entrenamiento está limitado por la carga de datos.
  \item \textbf{Máximo sobre las épocas.} Tomar el mejor valor de las tres épocas, y no el último, hace la medida robusta frente a la oscilación de una época concreta.
  \item \textbf{Números aleatorios comunes.} Cada evaluación restablece la semilla 42, de modo que dos individuos con la misma arquitectura parten de la misma inicialización y reciben los mismos minilotes en el mismo orden con las mismas transformaciones. Las diferencias de aptitud reflejan así los hiperparámetros y no el azar. Es una técnica clásica de reducción de varianza que hace más justas las comparaciones de la selección.
  \item \textbf{Tratamiento de restricciones.} Las configuraciones no factibles (divergencia numérica o falta de memoria) reciben aptitud nula, lo que equivale a una penalización de muerte. En la ejecución realizada ningún individuo resultó no factible.
  \item \textbf{Memoización.} Una tabla asocia cada fenotipo ya evaluado con su aptitud, de modo que un fenotipo repetido no consume presupuesto.
\end{itemize}

Esta aptitud tiene un sesgo intrínseco: premia las configuraciones que aprenden deprisa en las primeras épocas, que no tienen por qué ser las mejores tras un entrenamiento completo. Este sesgo, anticipado en el diseño, se analiza en la \cref{subsec:sesgo}.

\subsection{Operadores genéticos}\label{subsec:operadores}

\paragraph{Inicialización con siembra.} La población inicial de $N=10$ individuos contiene el genotipo de la configuración manual de referencia, $(0.500, 0.667, 0.500, 0.500)$, y nueve genotipos muestreados de forma uniforme en $[0,1]^4$. La siembra, combinada con el elitismo, garantiza que la solución final nunca sea peor, en aptitud, que la de referencia.

\paragraph{Selección por torneo.} Cada progenitor se elige tomando $k=3$ individuos al azar sin reemplazo y quedándose con el de mayor aptitud. Si los individuos se ordenan de mejor ($r=1$) a peor ($r=N$), la probabilidad de que gane el de rango $r$ es
\begin{equation}\label{eq:torneo}
  P(r) = \binom{N-r}{k-1}\bigg/\binom{N}{k},
\end{equation}
que para $N=10$ y $k=3$ vale 0.300, 0.233, 0.175, 0.125, 0.083, 0.050, 0.025 y 0.008 para los rangos 1 a 8, y cero para los dos peores. La presión selectiva es, por tanto, moderada: el mejor individuo tiene una probabilidad 1.3 veces mayor que el segundo y los dos peores nunca se reproducen. Además, la selección solo depende del orden y no de los valores de aptitud, lo que la hace invariante ante cambios de escala y robusta frente a diferencias pequeñas y ruidosas.

\paragraph{Cruce BLX-$\alpha$.} Con probabilidad $p_c=0.8$, cada pareja de progenitores $\vu^{a},\vu^{b}$ genera dos descendientes muestreando cada gen de forma independiente en un intervalo que extiende el segmento entre los progenitores~\cite{eshelman1993realcoded}:
\begin{equation}\label{eq:blx}
  u_j \sim U\bigl[\,m_j - \alpha\,d_j,\; M_j + \alpha\,d_j\,\bigr], \qquad m_j=\min(u^{a}_j,u^{b}_j),\; M_j=\max(u^{a}_j,u^{b}_j),\; d_j = M_j - m_j,
\end{equation}
con $\alpha=0.5$. El intervalo mide $(1+2\alpha)d_j=2d_j$: la mitad de la probabilidad cae entre los progenitores (explotación) y la otra mitad fuera de ellos (exploración por extrapolación). Este equilibrio evita la contracción sistemática de la población que provocaría un cruce puramente interpolante. Con probabilidad $1-p_c=0.2$, los descendientes son copias de los progenitores, que solo se diferenciarán por mutación.

\paragraph{Mutación gaussiana.} Cada gen de cada descendiente se perturba de forma independiente con probabilidad $p_m=0.25$:
\begin{equation}\label{eq:mutacion}
  u_j \leftarrow \rho\bigl(u_j + \epsilon_j\bigr), \qquad \epsilon_j \sim \mathcal{N}(0,\sigma^2),\quad \sigma = 0.2,
\end{equation}
de modo que se muta en promedio un gen por descendiente. En el espacio del fenotipo, $\sigma=0.2$ equivale a una desviación de 0.4 órdenes de magnitud en $\eta$ y de 0.6 en $\lambda$. En un gen categórico situado en el centro de su tramo, la mutación cambia el valor de $F$ o $D$ con probabilidad $P(|\epsilon|>1/6)\approx0.40$.

\paragraph{Reflexión en la frontera.} Los valores que salen de $[0,1]$ tras el cruce o la mutación se devuelven al intervalo por reflexión, $\rho(v) = 2-(|v| \bmod 2)$ si $|v|\bmod 2>1$ y $\rho(v)=|v|\bmod 2$ en otro caso. A diferencia del truncamiento, que acumularía probabilidad exactamente en los extremos ($\eta=10^{-4}$, $\lambda=10^{-6}$), la reflexión no introduce sesgo hacia la frontera.

\paragraph{Control de novedad.} Si el fenotipo de un descendiente ya ha sido evaluado o ya ha sido generado en la misma generación, se vuelve a mutar con $p_m=1$ hasta obtener un fenotipo inédito, con un máximo de 20 intentos. Este mecanismo, análogo a una memoria tabú, garantiza que cada una de las evaluaciones del presupuesto explore una configuración distinta y mantiene la diversidad de la población.

\paragraph{Reemplazo generacional con elitismo.} La nueva población está formada por el mejor individuo de la anterior ($e=1$), que no se reevalúa, y por $N-e=9$ descendientes. La mejor aptitud encontrada es, por tanto, monótona no decreciente.

\paragraph{Presupuesto.} Con $G=5$ generaciones, el número máximo de entrenamientos es $N+(G-1)(N-e)=10+4\cdot9=46$. La \cref{tab:parametros-ag} resume la configuración, la \cref{fig:flujo-ag} muestra el flujo del algoritmo y el \cref{alg:ag} lo describe en pseudocódigo.

\begin{table}[htbp]
  \centering
  \caption{Parámetros del algoritmo genético.}
  \label{tab:parametros-ag}
  \begin{tabular}{llr}
    \toprule
    \textbf{Componente} & \textbf{Parámetro} & \textbf{Valor} \\
    \midrule
    Población        & Tamaño $N$ / generaciones $G$ & 10 / 5 \\
    Selección        & Tamaño del torneo $k$ & 3 \\
    Cruce            & BLX-$\alpha$: probabilidad $p_c$ / $\alpha$ & 0.8 / 0.5 \\
    Mutación         & Gaussiana: probabilidad por gen $p_m$ / $\sigma$ & 0.25 / 0.2 \\
    Reemplazo        & Élite $e$ & 1 \\
    Presupuesto      & Evaluaciones máximas & 46 \\
    Aptitud          & Épocas / tamaño de lote / semilla & 3 / 8 / 42 \\
    \bottomrule
  \end{tabular}
\end{table}

\begin{figure}[htbp]
  \centering
  \begin{tikzpicture}[
      font=\small,
      paso/.style={draw=cazul!70!black, fill=fondogen, rounded corners=3pt, align=flush center, text width=5.4cm, minimum height=0.95cm, inner sep=4pt},
      oper/.style={draw=cnaranja!80!black, fill=fondocat, rounded corners=3pt, align=flush center, text width=5.8cm, minimum height=0.8cm, inner sep=4pt},
      fin/.style={draw=cverde!70!black, fill=cverde!12, rounded corners=3pt, align=flush center, text width=5.4cm, minimum height=0.8cm, inner sep=4pt},
      decision/.style={draw=cpurpura!80!black, fill=cpurpura!12, diamond, aspect=2.4, align=center, inner sep=1pt},
      flecha/.style={-{Stealth[length=2.2mm]}, line width=0.7pt, draw=black!70},
    ]
    \node[paso] (init) at (2.6, 0.0) {\textbf{Inicialización}\\referencia $+$ 9 genotipos aleatorios};
    \node[paso] (eval) at (2.6,-1.75) {\textbf{Evaluación de la aptitud}\\memoria de fenotipos o entrenamiento de 3 épocas};
    \node[decision] (dec) at (2.6,-3.6) {¿$g=G$?};
    \node[fin] (end) at (2.6,-5.35) {\textbf{Devolver el mejor individuo}};

    \node[oper] (elit) at (10.6,-5.35) {\textbf{Elitismo}: se conserva el mejor individuo};
    \node[oper] (sel)  at (10.4,-4.1)  {\textbf{Selección} por torneo, $k=3$};
    \node[oper] (cru)  at (10.4,-2.85) {\textbf{Cruce} BLX-$\alpha$, $p_c=0.8$, $\alpha=0.5$};
    \node[oper] (mut)  at (10.4,-1.6)  {\textbf{Mutación} gaussiana con reflexión\\$p_m=0.25$, $\sigma=0.2$};
    \node[oper] (nov)  at (10.4,-0.2)  {\textbf{Control de novedad}:\\remutar fenotipos ya evaluados};

    \draw[flecha] (init) -- (eval);
    \draw[flecha] (eval) -- (dec);
    \draw[flecha] (dec) -- node[right, font=\scriptsize] {sí} (end);
    \draw[flecha] (dec.east) -- node[above, font=\scriptsize] {no} (6.05,-3.6) |- (elit.west);
    \draw[flecha] (elit) -- (sel);
    \draw[flecha] (sel) -- (cru);
    \draw[flecha] (cru) -- (mut);
    \draw[flecha] (mut) -- (nov);
    \draw[flecha] (nov.west) -- (6.75,-0.2) |- node[pos=0.12, left, font=\scriptsize] {$g \leftarrow g+1$} (eval.east);

    \begin{scope}[on background layer]
      \node[draw=cnaranja!60, dashed, rounded corners=4pt, fit=(elit)(sel)(cru)(mut)(nov), inner sep=6pt, label={[font=\scriptsize, text=cnaranja!60!black]above:operadores de variación}] {};
    \end{scope}
  \end{tikzpicture}
  \caption{Flujo del algoritmo genético. La rama derecha genera los nueve descendientes de cada generación, que se suman al individuo élite para formar la siguiente población.}
  \label{fig:flujo-ag}
\end{figure}

\begin{algorithm}[htbp]
  \caption{Algoritmo genético para la optimización de hiperparámetros.}
  \label{alg:ag}
  \DontPrintSemicolon
  \KwIn{$N=10$, $G=5$, $e=1$, $k=3$, $p_c=0.8$, $\alpha=0.5$, $p_m=0.25$, $\sigma=0.2$; genotipo de referencia $\vu_{\mathrm{ref}}$; aptitud $\aptitud$; decodificación $\varphi$}
  \KwOut{mejor genotipo encontrado $\vu^{\star}$}
  $\mathcal{M} \leftarrow \emptyset$; \quad $P \leftarrow \{\vu_{\mathrm{ref}}\} \cup \{\vu_i \sim U[0,1]^4\}_{i=1}^{N-1}$\;
  \For{$g \leftarrow 1$ \KwTo $G$}{
    \ForEach{$\vu \in P$ sin aptitud asignada}{
      \eIf{$\varphi(\vu) \in \mathrm{dom}(\mathcal{M})$}{
        $f(\vu) \leftarrow \mathcal{M}[\varphi(\vu)]$\;
      }{
        $f(\vu) \leftarrow \aptitud(\vu)$ entrenando 3 épocas; \quad $\mathcal{M}[\varphi(\vu)] \leftarrow f(\vu)$\;
      }
    }
    ordenar $P$ por aptitud decreciente y actualizar $\vu^{\star}$\;
    \lIf{$g = G$}{\KwRet{$\vu^{\star}$}}
    $Q \leftarrow$ los $e$ mejores individuos de $P$; \quad $\mathcal{R} \leftarrow \mathrm{dom}(\mathcal{M})$\;
    \While{$|Q| < N$}{
      $\vu^{a} \leftarrow \textsc{Torneo}(P,k)$; \quad $\vu^{b} \leftarrow \textsc{Torneo}(P,k)$\;
      \eIf{$r \sim U[0,1]$ cumple $r < p_c$}{
        $\vu^{1} \leftarrow \textsc{BLX}_{\alpha}(\vu^{a},\vu^{b})$; \quad $\vu^{2} \leftarrow \textsc{BLX}_{\alpha}(\vu^{a},\vu^{b})$\;
      }{
        $\vu^{1} \leftarrow \vu^{a}$; \quad $\vu^{2} \leftarrow \vu^{b}$\;
      }
      \ForEach{$\vu \in \{\vu^{1},\vu^{2}\}$}{
        $\vu \leftarrow \textsc{Mutación}(\vu, p_m, \sigma)$\;
        \While{$\varphi(\vu) \in \mathcal{R}$ y no se hayan agotado 20 intentos}{
          $\vu \leftarrow \textsc{Mutación}(\vu, 1, \sigma)$\;
        }
        $\mathcal{R} \leftarrow \mathcal{R} \cup \{\varphi(\vu)\}$\;
        \lIf{$|Q| < N$}{$Q \leftarrow Q \cup \{\vu\}$}
      }
    }
    $P \leftarrow Q$\;
  }
\end{algorithm}

\subsection{Resultados de la búsqueda}\label{subsec:resultados-ag}

La búsqueda completó las 46 evaluaciones en 2\,h 17\,min de GPU. La \cref{tab:generaciones} resume cada generación y la \cref{fig:ag-evolucion} muestra la aptitud de todos los individuos evaluados.

\begin{table}[htbp]
  \centering
  \caption{Evolución de la población por generación. La media y la mediana incluyen al individuo élite.}
  \label{tab:generaciones}
  \small
  \begin{tabular}{crrrrrr}
    \toprule
    \textbf{Generación} & \textbf{Mejor} & \textbf{Media} & \textbf{Mediana} & \textbf{Mejor global} & \textbf{Evaluaciones} & \textbf{Minutos} \\
    \midrule
    1 & 0.4409 & 0.2464 & 0.2244 & 0.4409 & 10 & 27.2 \\
    2 & 0.4409 & 0.2983 & 0.3018 & 0.4409 & 19 & 52.1 \\
    3 & 0.5085 & 0.3327 & 0.3264 & 0.5085 & 28 & 80.9 \\
    4 & 0.5085 & 0.3560 & 0.3395 & 0.5085 & 37 & 109.2 \\
    5 & 0.5085 & 0.3370 & 0.2933 & 0.5085 & 46 & 136.8 \\
    \bottomrule
  \end{tabular}
\end{table}

\begin{figure}[htbp]
  \centering
  \begin{tikzpicture}
    \begin{axis}[
        informe,
        height=6.8cm,
        xlabel={Generación},
        ylabel={Aptitud (Dice de validación, 3 épocas)},
        xmin=0.5, xmax=5.5, ymin=0.08, ymax=0.56,
        xtick={1,2,3,4,5},
        legend style={at={(0.5,-0.2)}, anchor=north},
        legend columns=3,
        scatter/classes={
          8={mark=triangle*, mark size=2.6pt, cpurpura, draw=cpurpura!70!black},
          16={mark=*, mark size=2.2pt, cazul, draw=cazul!70!black},
          32={mark=square*, mark size=2.1pt, cverde, draw=cverde!70!black}
        },
      ]
      \addplot[scatter, only marks, scatter src=explicit symbolic, forget plot]
        table[col sep=comma,
              x expr={\thisrow{generacion} + (\thisrow{filtros}==8 ? -0.17 : (\thisrow{filtros}==32 ? 0.17 : 0))},
              y=aptitud, meta=filtros] {data/ga_individuos.csv};
      \addplot[only marks, mark=triangle*, mark size=2.6pt, cpurpura, draw=cpurpura!70!black] coordinates {(-1,0)};
      \addlegendentry{$F=8$}
      \addplot[only marks, mark=*, mark size=2.2pt, cazul, draw=cazul!70!black] coordinates {(-1,0)};
      \addlegendentry{$F=16$}
      \addplot[only marks, mark=square*, mark size=2.1pt, cverde, draw=cverde!70!black] coordinates {(-1,0)};
      \addlegendentry{$F=32$}
      \addplot[black, line width=1.4pt, const plot, mark=none] coordinates {(0.5,0.4409) (1,0.4409) (2,0.4409) (3,0.5085) (4,0.5085) (5,0.5085) (5.5,0.5085)};
      \addlegendentry{Mejor global}
      \addplot[cgris, line width=1pt, densely dashed, mark=diamond*, mark size=2.4pt, mark options={solid}] table[col sep=comma, x=generacion, y=media] {data/ga_generaciones.csv};
      \addlegendentry{Media de la población}
      \addplot[cbermellon, line width=0.9pt, dotted, domain=0.5:5.5, samples=2] {0.4409};
      \addlegendentry{Configuración de referencia}
    \end{axis}
  \end{tikzpicture}
  \caption{Aptitud de los 46 individuos evaluados, agrupados por generación y desplazados horizontalmente según el número de filtros base. La línea negra muestra la mejor aptitud encontrada hasta cada generación y la gris, la aptitud media de la población.}
  \label{fig:ag-evolucion}
\end{figure}

\paragraph{Dinámica de la búsqueda.} Durante las dos primeras generaciones, el mejor individuo es la configuración de referencia sembrada (0.4409). Sin embargo, la población mejora en promedio (la media pasa de 0.246 a 0.298), y en la segunda generación aparece un individuo con $\eta=3.16\cdot10^{-4}$, $\lambda=1.23\cdot10^{-6}$, 32 filtros y profundidad 4 que casi la iguala (0.4405). En la tercera generación surge la mejor solución, con aptitud 0.5085: $\eta=3.15\cdot10^{-4}$, $\lambda=1.91\cdot10^{-6}$, $F=16$ y $D=4$. Su genotipo es compatible con un descendiente de ese individuo de 32 filtros y de la configuración de referencia: conserva la tasa de aprendizaje y el decaimiento casi nulo del primero y la arquitectura de la segunda. El registro no almacena la genealogía, de modo que esta lectura es una hipótesis plausible y no un hecho verificado; en cualquier caso, ilustra cómo el cruce puede combinar bloques constructivos procedentes de individuos distintos. Las dos últimas generaciones no superan esa solución, aunque la cuarta alcanza la mayor aptitud media (0.356) y la quinta encuentra un individuo vecino con 0.4840.

\paragraph{Convergencia de la población.} La \cref{tab:convergencia-poblacion} muestra cómo la población se concentra en la región de la mejor solución. La proporción de individuos con 16 filtros pasa de 4 a 7 de cada 10, la de profundidad 4 de 4 a 8, y la mediana del decaimiento de pesos cae casi dos órdenes de magnitud. En la quinta generación ya no queda ningún individuo de 32 filtros. Es el comportamiento esperado de un AG: la selección explota la región prometedora mientras la mutación mantiene una exploración residual, visible en la caída de la aptitud media de la última generación.

\paragraph{Asignación del presupuesto frente al muestreo uniforme.} Esta concentración es la que justifica la elección del AG frente a la búsqueda aleatoria (\cref{subsubsec:aleatoria}). Con muestreo uniforme, cabría esperar que unas 5 de las 46 evaluaciones correspondieran a la arquitectura $F=16$, $D=4$ (una de las nueve combinaciones posibles), que unas 16 tuvieran una tasa de aprendizaje en la banda $[3\cdot10^{-4},\,1.5\cdot10^{-3}]$, en la que se encuentran todas las aptitudes superiores a 0.40, y que menos de 2 cumplieran ambas condiciones. El AG dedicó 18 evaluaciones a esa arquitectura, 23 a esa banda y 9 a su intersección. En las generaciones 3 a 5, casi la mitad de los individuos evaluados (13 de 27) tenían la arquitectura ganadora y tres de cada cuatro (20 de 27) un decaimiento $\lambda<10^{-5}$, frente al 11\,\% y el 33\,\% esperables por azar. El AG empleó, por tanto, su memoria para invertir el presupuesto en la región prometedora, algo que un muestreo sin memoria no puede hacer.

\begin{table}[htbp]
  \centering
  \caption{Composición de la población por generación (individuos de 10, incluido el élite) y medianas de los genes continuos.}
  \label{tab:convergencia-poblacion}
  \begin{tabular}{ccccc}
    \toprule
    \textbf{Generación} & \textbf{Con $F=16$} & \textbf{Con $D=4$} & \textbf{Mediana de $\eta$} & \textbf{Mediana de $\lambda$} \\
    \midrule
    1 & 4 & 4 & $1.54\cdot10^{-3}$ & $6.17\cdot10^{-5}$ \\
    2 & 4 & 5 & $6.75\cdot10^{-4}$ & $1.00\cdot10^{-4}$ \\
    3 & 7 & 8 & $7.63\cdot10^{-4}$ & $3.31\cdot10^{-6}$ \\
    4 & 7 & 8 & $5.32\cdot10^{-4}$ & $2.38\cdot10^{-6}$ \\
    5 & 6 & 7 & $3.84\cdot10^{-4}$ & $3.98\cdot10^{-6}$ \\
    \bottomrule
  \end{tabular}
\end{table}

\paragraph{Paisaje de la aptitud.} La \cref{fig:ag-paisaje} proyecta las 46 evaluaciones sobre los dos genes continuos, y la \cref{tab:marginales} resume la aptitud marginal de los genes categóricos. Las siete configuraciones con aptitud superior a 0.40 tienen una tasa de aprendizaje entre $3\cdot10^{-4}$ y $1.5\cdot10^{-3}$, mientras que por encima de $2\cdot10^{-3}$ ningún individuo supera 0.30. El decaimiento de pesos favorece valores bajos: seis de esas siete configuraciones tienen $\lambda<10^{-5}$, y la séptima es la propia referencia ($\lambda=10^{-4}$). En cuanto a la arquitectura, $F=16$ y $D=4$ obtienen la mayor aptitud media y máxima. Los extremos del espacio rinden peor en tres épocas: las redes pequeñas, por falta de capacidad, y las grandes, porque convergen más despacio.

\begin{figure}[htbp]
  \centering
  \begin{tikzpicture}
    \begin{groupplot}[
        group style={group size=2 by 1, horizontal sep=1.0cm, y descriptions at=edge left},
        informe,
        width=0.5\linewidth,
        height=5.6cm,
        ymin=0.08, ymax=0.56,
        xmode=log,
        scatter/classes={
          8={mark=triangle*, mark size=2.6pt, cpurpura, draw=cpurpura!70!black},
          16={mark=*, mark size=2.2pt, cazul, draw=cazul!70!black},
          32={mark=square*, mark size=2.1pt, cverde, draw=cverde!70!black}
        },
      ]
      \nextgroupplot[xlabel={Tasa de aprendizaje $\eta$}, ylabel={Aptitud}, xmin=8e-5, xmax=1.25e-2, legend to name=leyendapaisaje, legend columns=3]
      \addplot[scatter, only marks, scatter src=explicit symbolic] table[col sep=comma, x=lr, y=aptitud, meta=filtros] {data/ga_individuos.csv};
      \legend{$F=8$, $F=16$, $F=32$}
      \nextgroupplot[xlabel={Decaimiento de pesos $\lambda$}, xmin=8e-7, xmax=1.25e-3]
      \addplot[scatter, only marks, scatter src=explicit symbolic] table[col sep=comma, x=wd, y=aptitud, meta=filtros] {data/ga_individuos.csv};
    \end{groupplot}
  \end{tikzpicture}

  \vspace{2mm}
  \ref{leyendapaisaje}
  \caption{Paisaje de la aptitud observado: aptitud de los 46 individuos evaluados frente a la tasa de aprendizaje (izquierda) y al decaimiento de pesos (derecha), en escala logarítmica.}
  \label{fig:ag-paisaje}
\end{figure}

\begin{table}[htbp]
  \centering
  \caption{Aptitud marginal de los genes categóricos sobre las 46 evaluaciones.}
  \label{tab:marginales}
  \begin{tabular}{cccc@{\hspace{1.4cm}}cccc}
    \toprule
    $F$ & $n$ & \textbf{Media} & \textbf{Máx.} & $D$ & $n$ & \textbf{Media} & \textbf{Máx.} \\
    \midrule
    8  & 12 & 0.257 & 0.369 & 3 & 10 & 0.258 & 0.372 \\
    16 & 24 & \textbf{0.328} & \textbf{0.509} & 4 & 28 & \textbf{0.326} & \textbf{0.509} \\
    32 & 10 & 0.286 & 0.441 & 5 & 8  & 0.264 & 0.369 \\
    \bottomrule
  \end{tabular}
\end{table}

\paragraph{Interpretación de la solución.} El algoritmo genético conservó la arquitectura de la referencia ($F=16$, $D=4$), redujo la tasa de aprendizaje 3.2 veces y llevó el decaimiento de pesos a $1.91\cdot10^{-6}$, muy cerca del extremo inferior del dominio; en la práctica, desactivó la regularización $L_2$. Esta elección es coherente con la \cref{subsec:adam}: con Adam, la penalización acoplada es débil y poco uniforme, y en tres épocas cualquier regularización solo puede ralentizar el ajuste. La mejora de aptitud es de $+0.0676$, un 15.3\,\% relativo sobre la referencia.

\subsection{Validación de la solución: entrenamiento completo y conjunto de prueba}\label{subsec:validacion-ag}

La aptitud es una aproximación de tres épocas medida sobre el mismo conjunto de validación con el que se seleccionó la mejor de 46 configuraciones, de modo que su mejora es optimista por construcción. Para validar la solución, la configuración del AG (ejecución A) se reentrenó durante 100 épocas con un procedimiento idéntico al de la referencia, que se había entrenado durante 40, salvo en los hiperparámetros.

\begin{figure}[htbp]
  \centering
  \begin{tikzpicture}
    \begin{axis}[
        informe,
        height=6.6cm,
        xlabel={Época},
        ylabel={Dice medio de validación},
        xmin=0, xmax=100, ymin=0.1, ymax=0.85,
        legend pos=south east,
      ]
      \addplot[cazul!35, line width=0.5pt, forget plot] table[x=epoca, y=dice_val, col sep=comma] {data/historial_base.csv};
      \addplot[cbermellon!35, line width=0.5pt, forget plot] table[x=epoca, y=dice_val, col sep=comma] {data/historial_A.csv};
      \addplot[cazul, line width=1.3pt] table[x=epoca, y=dice_val_suave, col sep=comma] {data/historial_base.csv};
      \addlegendentry{Referencia manual}
      \addplot[cbermellon, line width=1.3pt] table[x=epoca, y=dice_val_suave, col sep=comma] {data/historial_A.csv};
      \addlegendentry{Configuración del AG (A)}
      \addplot[black!70, densely dashed, line width=0.8pt, domain=0:100, samples=2, forget plot] {0.6538};
      \addplot[only marks, mark=*, mark size=2.4pt, cazul, forget plot] coordinates {(35,0.6538)};
      \addplot[only marks, mark=*, mark size=2.4pt, cbermellon, forget plot] coordinates {(10,0.6663) (64,0.7702)};
      \draw[black!60, line width=0.4pt] (axis cs:35,0.6538) -- (axis cs:43,0.585);
      \node[font=\scriptsize, anchor=north west, align=left, fill=white, inner sep=1.5pt] at (axis cs:43,0.59) {mejor de la referencia:\\0.6538 (época 35, 32 min)};
      \draw[black!60, line width=0.4pt] (axis cs:10,0.6663) -- (axis cs:14,0.78);
      \node[font=\scriptsize, anchor=south west, align=left, fill=white, inner sep=1.5pt] at (axis cs:12,0.78) {A supera ese valor en la época 10 (10 min)};
      \node[font=\scriptsize, anchor=south west, fill=white, inner sep=1.5pt] at (axis cs:61,0.785) {mejor de A: 0.7702 (época 64)};
    \end{axis}
  \end{tikzpicture}
  \caption{Dice medio de validación por época. Las líneas finas muestran los valores por época y las gruesas, la media móvil centrada de 5 épocas.}
  \label{fig:validacion-ag}
\end{figure}

\paragraph{Velocidad de convergencia.} La \cref{fig:validacion-ag} muestra que la configuración del AG supera el mejor valor de validación de la referencia (0.6538, alcanzado en la época 35 tras 32 minutos) en la época 10, tras 10 minutos de entrenamiento: una convergencia 3.2 veces más rápida. Dentro del mismo presupuesto de 40 épocas, A alcanza 0.7354 frente a 0.6538. Su mejor valor, 0.7702, se obtiene en la época 64.

\paragraph{Precisión en el conjunto de prueba.} La \cref{tab:prueba-ag} recoge la comparación sobre los 592 cortes de prueba, nunca vistos durante la búsqueda ni la selección de puntos de control. El Dice medio pasa de 0.6905 a 0.7760 ($+0.0855$) y el IoU medio de 0.5519 a 0.6470. La mejora se concentra en el glioma, la clase más difícil ($+0.178$ de Dice), mientras que el tumor hipofisario apenas cambia.

\begin{table}[htbp]
  \centering
  \caption{Comparación en el conjunto de prueba (592 cortes, 46 pacientes) de la configuración manual de referencia y la configuración del algoritmo genético.}
  \label{tab:prueba-ag}
  \begin{tabular}{lrrrrrr}
    \toprule
    & \multicolumn{3}{c}{\textbf{Dice}} & \multicolumn{3}{c}{\textbf{IoU}} \\
    \cmidrule(lr){2-4}\cmidrule(lr){5-7}
    \textbf{Región} & \textbf{Referencia} & \textbf{AG (A)} & $\Delta$ & \textbf{Referencia} & \textbf{AG (A)} & $\Delta$ \\
    \midrule
    Meningioma   & 0.8258 & \textbf{0.8899} & $+0.0641$ & 0.7033 & \textbf{0.8017} & $+0.0984$ \\
    Glioma       & 0.4471 & \textbf{0.6254} & $+0.1784$ & 0.2879 & \textbf{0.4550} & $+0.1671$ \\
    Hipofisario  & 0.7985 & \textbf{0.8126} & $+0.0141$ & 0.6646 & \textbf{0.6844} & $+0.0198$ \\
    \midrule
    Media (3 clases) & 0.6905 & \textbf{0.7760} & $+0.0855$ & 0.5519 & \textbf{0.6470} & $+0.0951$ \\
    Tumor completo   & 0.6230 & \textbf{0.7233} & $+0.1003$ & 0.4525 & \textbf{0.5666} & $+0.1141$ \\
    \bottomrule
  \end{tabular}
\end{table}

\paragraph{Lectura crítica.} La mejora en prueba ($+0.0855$) es menor que en validación ($+0.1164$), como cabía esperar: parte de la ventaja de validación se debía al optimismo de seleccionar la mejor de 46 configuraciones sobre ese mismo conjunto. Además, los presupuestos de entrenamiento difieren (40 frente a 100 épocas), por lo que la comparación en prueba mezcla el efecto de los hiperparámetros con el de un entrenamiento más largo. La comparación con igual presupuesto en validación (0.7354 frente a 0.6538 en 40 épocas) indica que la mayor parte de la ganancia procede de los hiperparámetros. Por último, la evidencia disponible compara el AG con una configuración manual, no con otras estrategias de búsqueda. La elección del AG frente a la búsqueda aleatoria se apoya en los argumentos teóricos de la \cref{subsubsec:aleatoria} y en la asignación del presupuesto observada en la \cref{subsec:resultados-ag}, pero no en una comparación empírica directa. Se implementó una búsqueda aleatoria de referencia con el mismo espacio, la misma aptitud y el mismo presupuesto de 46 evaluaciones, pero no se dispone de sus resultados, por lo que esa comparación queda como trabajo futuro.

\section{Superando los límites de la heurística: análisis y refinamiento}\label{sec:refinamiento}

El algoritmo genético encontró una configuración del optimizador claramente mejor que la manual, pero su solución está condicionada por dos limitaciones de diseño: la aptitud de tres épocas y un cromosoma que deja fuera la función de pérdida, el aumento de datos y la programación de la tasa de aprendizaje. Esta sección describe cómo se analizó la solución heurística y cómo se superaron esas limitaciones mediante un ciclo de diagnóstico, formulación de hipótesis y experimentación controlada (Fases 5 y 6). En todos los experimentos se conservan la tasa de aprendizaje y el decaimiento de pesos encontrados por el AG.

\subsection{Diagnóstico de la solución del algoritmo genético}\label{subsec:diagnostico}

La \cref{tab:diagnostico-A} desglosa el error del modelo A en validación. El hallazgo principal es que el glioma está \emph{limitado por la sensibilidad}: el modelo omite el 44.2\,\% de sus píxeles, que clasifica como fondo, mientras que solo el 1.4\,\% se confunde con otro tipo de tumor. Su razón FN/FP de 2.25 indica que pierde algo más del doble de píxeles tumorales de los que marca por error. Además, en el 14.4\,\% de los cortes de glioma no se predice ningún píxel tumoral, y el Dice medio por corte de los gliomas pequeños (tercio inferior por área) es de solo 0.380, frente a 0.672 en los grandes. El problema no es reconocer el tipo de tumor, sino delimitar sus bordes difusos.

\begin{table}[htbp]
  \centering
  \caption{Diagnóstico por clase del modelo A (configuración del AG) en validación.}
  \label{tab:diagnostico-A}
  \small
  \begin{tabular}{lcccccc}
    \toprule
    \textbf{Clase} & \textbf{Precisión} & \textbf{Sensibilidad} & \textbf{Dice} & \textbf{FN/FP} & \textbf{Omitido} & \textbf{Confundido} \\
    & & & & & \textbf{como fondo} & \textbf{de tipo} \\
    \midrule
    Meningioma  & 0.907 & 0.821 & 0.862 & 2.13 & 11.6\,\% & 6.3\,\% \\
    Glioma      & 0.728 & 0.544 & 0.623 & 2.25 & 44.2\,\% & 1.4\,\% \\
    Hipofisario & 0.808 & 0.844 & 0.826 & 0.78 & 10.8\,\% & 4.8\,\% \\
    \bottomrule
  \end{tabular}
\end{table}

A la luz de la \cref{eq:gradiente-ti}, la corrección natural es una pérdida que premie más la detección de los píxeles tumorales que la evitación de falsos positivos: la pérdida Focal Tversky con $\alpha=0.7>\beta=0.3$.

\subsection{El sesgo de horizonte corto de la función de aptitud}\label{subsec:sesgo}

La aptitud del AG mide el Dice tras tres épocas, y por tanto premia la velocidad inicial de aprendizaje y no la calidad final. Este fenómeno, conocido como \emph{sesgo de horizonte corto} en la metaoptimización~\cite{wu2018shorthorizon}, se observa con claridad en los datos de la búsqueda:

\begin{itemize}
  \item \textbf{Penaliza la capacidad.} La aptitud media de los individuos de 32 filtros (0.286) es menor que la de los de 16 (0.328). El caso más elocuente es el individuo de la segunda generación con $\eta=3.16\cdot10^{-4}$, $\lambda=1.23\cdot10^{-6}$, $F=32$ y $D=4$: tiene prácticamente la misma configuración del optimizador que la solución ganadora y, sin embargo, obtiene 0.4405 frente a 0.5085. En tres épocas, una red cuatro veces mayor aún no ha avanzado tanto en su descenso, lo que no significa que su punto final vaya a ser peor.
  \item \textbf{Es ciega a la programación de la tasa de aprendizaje.} Con recocido coseno sobre $T=150$ épocas, la tasa tras tres épocas es el 99.9\,\% de la inicial (\cref{eq:coseno}). Cualquier beneficio de reducir el paso al final del entrenamiento es invisible para la aptitud.
  \item \textbf{Penaliza la regularización.} El aumento de datos y las pérdidas que reequilibran precisión y sensibilidad ralentizan el ajuste inicial, aunque mejoren la generalización. Además, ninguno de ellos formaba parte del cromosoma.
\end{itemize}

De este análisis se derivan cuatro hipótesis, que se contrastaron con experimentos controlados: (H1) la pérdida Focal Tversky mejora la sensibilidad en el glioma; (H2) el aumento de datos regulariza y mejora la generalización; (H3) la tasa de aprendizaje constante provoca oscilaciones que convierten la selección del punto de control en un sorteo; y (H4) la capacidad de 16 filtros, elegida bajo el sesgo de la aptitud, limita el rendimiento final.

\subsection{Fase 5: diseño factorial de la pérdida y el aumento de datos}\label{subsec:fase5}

Para contrastar H1 y H2 se entrenó un diseño factorial $2\times2$ (pérdida Dice-Focal o Focal Tversky-Focal, aumento básico o avanzado; ejecuciones A, B, C y D) y, a continuación, dos ejecuciones de aislamiento que separan los bloques espacial y de intensidad del aumento avanzado (S e I). La \cref{tab:fase5} resume los resultados en validación.

\begin{table}[htbp]
  \centering
  \caption{Fase 5: resultados en validación (mejor punto de control). DF: Dice-Focal; FT: Focal Tversky-Focal con $\alpha=0.7$, $\beta=0.3$, $\gamma_T=4/3$. \emph{Peq.}: Dice medio por corte de los gliomas pequeños. \emph{Omit.}: porcentaje de cortes de glioma sin ningún píxel tumoral predicho.}
  \label{tab:fase5}
  \small
  \setlength{\tabcolsep}{3.4pt}
  \begin{tabular}{llcccccccc}
    \toprule
    & & \multicolumn{4}{c}{\textbf{Dice}} & \multicolumn{4}{c}{\textbf{Glioma}} \\
    \cmidrule(lr){3-6}\cmidrule(lr){7-10}
    \textbf{Ejec.} & \textbf{Aumento y pérdida} & \textbf{Medio} & \textbf{Men.} & \textbf{Gli.} & \textbf{Hip.} & \textbf{Sens.} & \textbf{FN/FP} & \textbf{Peq.} & \textbf{Omit.} \\
    \midrule
    A & básico, DF            & 0.7702 & \textbf{0.862} & 0.623 & 0.826 & 0.544 & 2.25 & 0.380 & 14.4\,\% \\
    B & avanzado, DF          & 0.7595 & 0.852 & 0.629 & 0.797 & 0.584 & 1.54 & 0.390 & 14.4\,\% \\
    C & básico, FT            & 0.7672 & 0.856 & 0.622 & 0.824 & 0.582 & 1.44 & 0.312 & 12.2\,\% \\
    D & avanzado, FT          & 0.7675 & 0.845 & \textbf{0.683} & 0.774 & 0.662 & 1.22 & \textbf{0.426} & 10.1\,\% \\
    \midrule
    S & espacial avanzado, FT & \textbf{0.7751} & 0.858 & 0.640 & \textbf{0.828} & \textbf{0.694} & 0.64 & 0.404 & \textbf{8.0\,\%} \\
    I & intensidad, FT        & 0.7454 & 0.830 & 0.618 & 0.788 & 0.568 & 1.60 & 0.360 & 17.0\,\% \\
    \bottomrule
  \end{tabular}
\end{table}

Los resultados permiten extraer las siguientes conclusiones:

\begin{enumerate}
  \item \textbf{La pérdida Focal Tversky desplaza el equilibrio entre precisión y sensibilidad tal como predice el análisis del gradiente (H1).} Por sí sola (C frente a A), eleva la sensibilidad del glioma en 0.038 y reduce su razón FN/FP de 2.25 a 1.44. Sin embargo, su precisión cae en 0.061 y el Dice del glioma no varía. Además, la ejecución sobreajusta: la pérdida de validación alcanza su mínimo en la época 52 y después crece de 0.137 a 0.188.
  \item \textbf{El aumento espacial es el componente útil del aumento avanzado (H2).} Combinada con el aumento avanzado (D), la pérdida Focal Tversky produce la mayor ganancia en el glioma. El aislamiento atribuye esa ganancia al bloque espacial: frente a C, la ejecución S eleva la sensibilidad del glioma de 0.582 a 0.694, el Dice de los gliomas pequeños de 0.312 a 0.404 y reduce los cortes de glioma omitidos del 12.2\,\% al 8.0\,\%. El bloque de intensidad, en cambio, es perjudicial: la ejecución I omite más gliomas (17.0\,\%) y genera una confusión de meningioma con tumor hipofisario (el 5.9\,\% de los píxeles de meningioma) que no aparece en las demás.
  \item \textbf{Las diferencias en Dice medio están dentro del ruido.} La mejor configuración, S (aumento espacial avanzado con Focal Tversky), supera a A en solo $+0.005$, mientras que el Dice de validación oscila entre $0.02$ y $0.03$ de una época a la siguiente. Entre las épocas 21 y 42, ambas ejecuciones tienen exactamente la misma media (0.699). La Fase 5 cambió el perfil de error del glioma, pero no el rendimiento medio. Por la misma razón, la ganancia de D en el glioma en su mejor punto de control (época 97) sobrestima el efecto: promediando épocas comparables, el aumento aporta entre 0.02 y 0.04 puntos de Dice en el glioma.
  \item \textbf{Con aumento espacial, $\alpha=0.7$ sobrecorrige.} En S, la razón FN/FP del glioma cae a 0.64 y su precisión a 0.593: el modelo ha pasado a estar limitado por la precisión. Según la \cref{eq:gradiente-ti}, la corrección consiste en reducir el peso relativo de los falsos negativos, lo que lleva a $\alpha=0.6$ y $\beta=0.4$.
\end{enumerate}

En conjunto, siete configuraciones distintas alcanzan su mejor valor entre 0.744 y 0.775, y sus curvas suavizadas se estabilizan en torno a 0.70--0.74. Esta meseta, junto con la oscilación de una época a otra, apunta a H3 y H4: el límite no estaba en la pérdida ni en el aumento, sino en la dinámica del descenso y en la capacidad del modelo.

\subsection{Fase 6: capacidad y estabilización del descenso}\label{subsec:fase6}

La ejecución E combina las tres correcciones derivadas del análisis: 32 filtros base (7.24 millones de parámetros), recocido coseno de la tasa de aprendizaje sobre $T=150$ épocas y pesos Focal Tversky recalibrados ($\alpha=0.6$, $\beta=0.4$), junto con el aumento espacial avanzado y los valores de $\eta$ y $\lambda$ del AG. Como E modifica tres factores a la vez respecto a S, se entrenó también un control, E16, idéntico salvo por tener 16 filtros: la comparación E frente a E16 aísla la capacidad, y E16 frente a S aísla el recocido junto con los nuevos pesos.

\begin{figure}[htbp]
  \centering
  \begin{tikzpicture}
    \begin{groupplot}[
        group style={group size=1 by 2, vertical sep=0.35cm, x descriptions at=edge bottom},
        informe,
        xmin=0, xmax=150,
      ]
      \nextgroupplot[height=6.3cm, ylabel={Dice medio de validación}, ymin=0.25, ymax=0.82, legend pos=south east, legend columns=3]
      \addplot[cgris, line width=0.8pt] table[x=epoca, y=dice_val, col sep=comma] {data/historial_A.csv};
      \addlegendentry{A: constante, $F=16$}
      \addplot[cazul, line width=0.8pt] table[x=epoca, y=dice_val, col sep=comma] {data/historial_E16.csv};
      \addlegendentry{E16: coseno, $F=16$}
      \addplot[cbermellon, line width=0.9pt] table[x=epoca, y=dice_val, col sep=comma] {data/historial_E.csv};
      \addlegendentry{E: coseno, $F=32$}
      \nextgroupplot[height=3.4cm, xlabel={Época}, ylabel={$\eta_e\;(\times10^{-4})$}, ymin=0, ymax=3.6, ytick={0,1,2,3}]
      \addplot[cgris, line width=0.9pt, domain=0:100, samples=2] {3.15144};
      \addplot[cbermellon, line width=0.9pt] table[x=epoca, y expr={\thisrow{lr}*10000}, col sep=comma] {data/historial_E.csv};
    \end{groupplot}
  \end{tikzpicture}
  \caption{Arriba: Dice medio de validación por época (valores sin suavizar) con tasa de aprendizaje constante (A) y con recocido coseno (E16 y E). Abajo: tasa de aprendizaje aplicada en cada época. Al reducirse el paso, las oscilaciones desaparecen y el modelo se asienta.}
  \label{fig:coseno}
\end{figure}

\begin{table}[htbp]
  \centering
  \caption{Fase 6: dinámica del entrenamiento en validación. El cambio medio entre épocas consecutivas se calcula sobre las 30 últimas épocas de cada ejecución, y la desviación típica y el rango, sobre las 20 últimas.}
  \label{tab:dinamica}
  \small
  \begin{tabular}{lcccc}
    \toprule
    & \textbf{A} & \textbf{S} & \textbf{E16} & \textbf{E} \\
    \midrule
    Filtros base / programación de $\eta$ & 16 / const. & 16 / const. & 16 / coseno & 32 / coseno \\
    Mejor época / mejor Dice medio & 64 / 0.7702 & 51 / \textbf{0.7751} & 95 / 0.7693 & 72 / 0.7612 \\
    Cambio medio entre épocas & 0.022 & 0.030 & \textbf{0.003} & \textbf{0.003} \\
    Desviación típica & 0.016 & 0.019 & \textbf{0.002} & \textbf{0.002} \\
    Rango & 0.053 & 0.068 & 0.011 & \textbf{0.006} \\
    Mejor media móvil de 5 épocas & 0.749 & 0.735 & \textbf{0.756} & 0.755 \\
    Media de las 10 últimas épocas & 0.730 & 0.719 & 0.741 & \textbf{0.749} \\
    \bottomrule
  \end{tabular}
\end{table}

La \cref{fig:coseno} y la \cref{tab:dinamica} confirman las hipótesis H3 y H4, con matices:

\begin{enumerate}
  \item \textbf{El recocido coseno elimina las oscilaciones (H3).} El cambio medio del Dice de validación entre épocas consecutivas cae de 0.022--0.030 a 0.003, y su rango en las últimas 20 épocas de 0.05--0.07 a 0.006--0.011. La última época pasa a ser un modelo utilizable (0.749 en E), y la selección del punto de control deja de ser un sorteo. Es la predicción de la \cref{subsec:programacion}: con pasos decrecientes, los iterados dejan de oscilar alrededor del punto estacionario y se asientan en él.
  \item \textbf{La capacidad aporta una mejora pequeña pero consistente en convergencia (H4).} En las 30 épocas finales, ya recocidas, E supera a E16 en $+0.009$ de Dice medio, gracias sobre todo al glioma ($+0.036$) y al meningioma ($+0.025$), aunque el tumor hipofisario retrocede ($-0.035$). En el mejor punto de control, en cambio, E queda por debajo de E16 (0.7612 frente a 0.7693), porque ambos picos caen en la fase intermedia, todavía ruidosa, del recocido. Gran parte de la capacidad adicional se dedica a ajustar el conjunto de entrenamiento: la pérdida de entrenamiento final es de $\approx0.02$ en ambos casos. Como el entrenamiento estaba limitado por la carga de datos, cuadruplicar los parámetros no aumentó el tiempo por época (31.2\,s en E frente a 33.3\,s en E16).
  \item \textbf{Los pesos $\alpha=0.6$, $\beta=0.4$ equilibran el glioma, pero solo con 32 filtros.} En E, la precisión del glioma sube de 0.593 a 0.702 respecto a S, su sensibilidad apenas varía (de 0.694 a 0.679) y la razón FN/FP queda en 1.11, prácticamente equilibrada. Con los mismos pesos, E16 vuelve al perfil inicial (FN/FP de 2.10). El equilibrio depende de la combinación de pérdida y capacidad, no de $\alpha$ y $\beta$ por separado: los hiperparámetros interactúan, como se anticipó en la \cref{subsec:caracterizacion}.
\end{enumerate}

\paragraph{Resultado en el conjunto de prueba.} La comparación planificada de A frente a E en prueba (\cref{tab:prueba-E}) muestra una mejora de $+0.0412$ en Dice medio, hasta \textbf{0.8175}, el mejor resultado del proyecto con una red entrenada desde cero. La mejora se concentra de nuevo en el glioma ($+0.0876$), seguido del tumor hipofisario ($+0.0328$). Respecto a la configuración manual de referencia, la ganancia total es de $+0.127$ (un 18.4\,\% relativo), y el Dice del glioma pasa de 0.447 a 0.714.

\begin{table}[htbp]
  \centering
  \caption{Conjunto de prueba (592 cortes): referencia manual, configuración del AG (A) y modelo refinado (E). La evaluación de A se repitió en esta fase con un resultado prácticamente idéntico al de la \cref{tab:prueba-ag} (0.7763 frente a 0.7760).}
  \label{tab:prueba-E}
  \begin{tabular}{lcccc}
    \toprule
    \textbf{Región} & \textbf{Referencia} & \textbf{AG (A)} & \textbf{Refinado (E)} & \textbf{E $-$ A} \\
    \midrule
    Dice meningioma  & 0.8258 & 0.8901 & \textbf{0.8935} & $+0.0034$ \\
    Dice glioma      & 0.4471 & 0.6259 & \textbf{0.7135} & $+0.0876$ \\
    Dice hipofisario & 0.7985 & 0.8129 & \textbf{0.8457} & $+0.0328$ \\
    \midrule
    \textbf{Dice medio} & 0.6905 & 0.7763 & \textbf{0.8175} & $+0.0412$ \\
    IoU medio        & 0.5519 & 0.6474 & \textbf{0.6982} & $+0.0508$ \\
    Dice tumor completo & 0.6230 & 0.7237 & \textbf{0.7762} & $+0.0525$ \\
    \bottomrule
  \end{tabular}
\end{table}

\paragraph{Discrepancia entre validación y prueba.} En validación, E queda por debajo de A (0.7612 frente a 0.7702); en prueba, claramente por encima (0.8175 frente a 0.7763). La mejora del glioma es el único efecto que se replica en ambos conjuntos ($+0.068$ en validación y $+0.088$ en prueba) y, por tanto, el más creíble. Las diferencias en meningioma y tumor hipofisario cambian de signo de un conjunto a otro, lo que es coherente con la inestabilidad entre puntos de control observada en la Fase 5. La conclusión metodológica es que un conjunto de validación de 468 cortes no permite ordenar modelos con una resolución de $\pm0.01$--$0.02$. Esta limitación afecta también a la aptitud del AG, que se mide sobre ese mismo conjunto.

\subsection{Síntesis: búsqueda global heurística y refinamiento guiado}\label{subsec:sintesis}

El recorrido de las Fases 3 a 6 puede leerse como una búsqueda híbrida en dos etapas, análoga a un algoritmo memético. El algoritmo genético realizó una exploración global, barata y automática del espacio del optimizador y de la arquitectura, y encontró una cuenca buena: la tasa de aprendizaje y el decaimiento que se mantuvieron hasta el modelo final. El refinamiento posterior actuó como una búsqueda local guiada por diagnóstico, en dimensiones que el cromosoma no contemplaba (pérdida, aumento y programación de la tasa) y corrigiendo el sesgo de la aptitud (capacidad). Sobre el conjunto de prueba, la primera etapa aporta $+0.086$ puntos de Dice medio (unos dos tercios de la mejora) y la segunda $+0.041$ (un tercio).

\section{Complemento: transferencia de aprendizaje}\label{sec:transferencia}

Como experimento adicional, ajeno al núcleo de optimización del trabajo, se evaluó si un extractor de características externo puede superar los límites intrínsecos de un conjunto de datos pequeño. La ejecución F sustituye el codificador entrenado desde cero por una ResNet34~\cite{he2016resnet} preentrenada en ImageNet~\cite{deng2009imagenet}, integrada en una U-Net de la biblioteca Segmentation Models PyTorch~\cite{iakubovskii2019smp} con un único canal de entrada y cuatro clases de salida (24.43 millones de parámetros). El resto de la configuración se hereda de E: aumento espacial avanzado, pérdida Focal Tversky-Focal con $\alpha=0.6$ y $\beta=0.4$, recocido coseno (100 épocas) y la tasa de aprendizaje y el decaimiento del AG. La única adaptación es una tasa de aprendizaje diferenciada: el decodificador usa la tasa del AG ($3.15\cdot10^{-4}$) y el codificador preentrenado una diez veces menor, para que los gradientes grandes y ruidosos del decodificador recién inicializado no destruyan las características de ImageNet en las primeras épocas.

\begin{table}[htbp]
  \centering
  \caption{Modelo refinado entrenado desde cero (E) frente al modelo con codificador preentrenado (F). Los intervalos son percentiles al 95\,\% sobre 10\,000 remuestreos pareados de los 46 pacientes de prueba, estratificados por tipo de tumor.}
  \label{tab:transferencia}
  \small
  \begin{tabular}{lcccc}
    \toprule
    & \textbf{E (desde cero)} & \textbf{F (ResNet34)} & \textbf{F $-$ E} & \textbf{IC 95\,\%} \\
    \midrule
    Dice medio en validación  & 0.7612 & \textbf{0.8221} & $+0.061$ & -- \\
    \midrule
    Dice meningioma (prueba)  & 0.8935 & \textbf{0.9221} & $+0.029$ & $[+0.008,\,+0.074]$ \\
    Dice glioma (prueba)      & \textbf{0.7135} & 0.6855 & $-0.028$ & $[-0.066,\,+0.004]$ \\
    Dice hipofisario (prueba) & 0.8457 & \textbf{0.8672} & $+0.022$ & $[-0.004,\,+0.046]$ \\
    \textbf{Dice medio (prueba)} & 0.8175 & \textbf{0.8250} & $+0.007$ & $[-0.010,\,+0.029]$ \\
    \midrule
    Épocas hasta Dice val. $\geq0.75$ & 49 & \textbf{7} & & \\
    Tiempo de entrenamiento & 78 min & \textbf{47 min} & & \\
    \bottomrule
  \end{tabular}
\end{table}

La \cref{tab:transferencia} resume los resultados, que admiten tres lecturas:

\begin{itemize}
  \item \textbf{En Dice medio, E y F están estadísticamente empatados.} F obtiene 0.8250 en prueba (intervalo de confianza al 95\,\% de $[0.785,\,0.852]$), pero su ventaja sobre E, $+0.007$, tiene un intervalo que incluye el cero. La gran ventaja en validación ($+0.061$) no se traslada a prueba, lo que confirma de nuevo que el conjunto de validación es un instrumento optimista.
  \item \textbf{Las características externas resuelven el reconocimiento del tipo de tumor.} La confusión de tumor hipofisario con meningioma en validación cae del 18.4\,\% al 4.2\,\%, y en prueba ninguna confusión entre tipos supera el 1\,\% de los píxeles. La mejora en meningioma es significativa. Es el beneficio más claro de ImageNet: un contexto global más rico para distinguir lesiones de apariencia parecida.
  \item \textbf{No mejoran la delimitación del glioma.} F es más conservador: su precisión en el glioma sube (0.894 en prueba) pero su sensibilidad cae (0.556), con una razón FN/FP de 6.7, y empeora en 12 de los 18 pacientes con glioma. Los bordes infiltrativos siguen siendo el límite, posiblemente porque un único corte T1-CE ofrece una señal débil para delimitarlos.
\end{itemize}

La ventaja inequívoca es la eficiencia: F alcanza un Dice de validación de 0.75 en la época 7, frente a la 49 de E, y termina su entrenamiento en 47 minutos frente a 78. La transferencia de aprendizaje acelera de forma drástica el descenso, porque parte de un punto inicial ya cercano a una buena cuenca, pero en este conjunto de datos no desplaza el punto final de forma significativa. Este resultado refuerza la conclusión principal del trabajo: con una arquitectura entrenada desde cero, la combinación de optimización heurística de hiperparámetros y refinamiento guiado alcanza un rendimiento equivalente al de un codificador preentrenado con más del triple de parámetros.

\section{Discusión y conclusiones}\label{sec:discusion}

\subsection{Dos niveles de optimización complementarios}\label{subsec:dos-niveles}

El trabajo ha resuelto dos problemas de optimización de naturaleza opuesta, cuyas características se contrastan en la \cref{tab:dos-niveles}. El nivel interno es un problema de dimensión enorme con información de primer orden exacta y barata, que un método de descenso local resuelve eficazmente. El nivel externo es un problema de dimensión mínima sin ninguna información de gradiente, en el que cada evaluación es carísima, y que solo una metaheurística poblacional puede abordar con un presupuesto de unas decenas de evaluaciones. Ninguno de los dos métodos serviría para el problema del otro: un algoritmo genético no podría ajustar millones de pesos, y el descenso de gradiente no puede derivar el Dice de validación respecto a la profundidad de la red.

\begin{table}[htbp]
  \centering
  \caption{Comparación de los dos niveles de optimización del trabajo.}
  \label{tab:dos-niveles}
  \small
  \begin{tabularx}{\linewidth}{>{\raggedright\arraybackslash}p{3.1cm}XX}
    \toprule
    & \textbf{Nivel interno (U-Net)} & \textbf{Nivel externo (algoritmo genético)} \\
    \midrule
    Variables de decisión & Pesos $\vtheta\in\R^{d}$, con $d\approx1.8\cdot10^{6}$--$7.2\cdot10^{6}$ & Genotipo $\vu\in[0,1]^4$ que codifica $(\eta,\lambda,F,D)$ \\
    Función objetivo & Pérdida sustituta diferenciable (Dice-Focal o Focal Tversky-Focal) & Dice de validación tras 3 épocas (aproximación de baja fidelidad) \\
    Información disponible & Gradiente exacto por retropropagación & Solo el valor de la aptitud (caja negra) \\
    Método & Adam: descenso estocástico adaptativo de primer orden & Torneo, cruce BLX-$\alpha$, mutación gaussiana y elitismo \\
    Coste por evaluación & Una iteración: milisegundos & Un entrenamiento de 3 épocas: unos 3 minutos \\
    Naturaleza de la búsqueda & Local, desde una inicialización & Global y poblacional \\
    Resultado & Pesos de cada modelo & $\eta=3.15\cdot10^{-4}$, $\lambda=1.91\cdot10^{-6}$, $F=16$, $D=4$ \\
    \bottomrule
  \end{tabularx}
\end{table}

Los dos niveles están acoplados en ambos sentidos. Hacia abajo, los hiperparámetros elegidos por el AG determinan la trayectoria del descenso: la \cref{fig:perdidas} muestra que la misma red con la misma pérdida desciende más deprisa y llega más lejos con la configuración evolucionada. Hacia arriba, las propiedades del descenso condicionan lo que el AG puede ver: como la aptitud trunca el entrenamiento a tres épocas, el AG solo percibe la fase inicial de la trayectoria, y de ahí su sesgo hacia redes pequeñas y su ceguera ante la programación de la tasa de aprendizaje.

\subsection{Contribución de cada etapa}\label{subsec:contribucion}

La \cref{fig:prueba-final} y la \cref{tab:contribucion} resumen la evolución del rendimiento en el conjunto de prueba. La metaoptimización con el algoritmo genético es, con diferencia, la etapa que más aporta: el 63.8\,\% de la mejora total y el 67.6\,\% de la obtenida con redes entrenadas desde cero. El refinamiento guiado por diagnóstico aporta casi un tercio, y la transferencia de aprendizaje, una mejora pequeña y no significativa.

\begin{figure}[htbp]
  \centering
  \begin{tikzpicture}
    \begin{axis}[
        informe,
        ybar,
        height=6.6cm,
        bar width=9pt,
        enlarge x limits=0.15,
        ymin=0, ymax=1.0,
        ylabel={Dice en el conjunto de prueba},
        symbolic x coords={Meningioma,Glioma,Hipofisario,Media},
        xtick=data,
        xticklabel style={font=\small},
        legend style={at={(0.5,1.03)}, anchor=south, /tikz/every even column/.append style={column sep=8pt}},
        legend columns=4,
        grid=major,
        xmajorgrids=false,
      ]
      \addplot[fill=cgris!55, draw=cgris!80!black] coordinates {(Meningioma,0.8258) (Glioma,0.4471) (Hipofisario,0.7985) (Media,0.6905)};
      \addlegendentry{Referencia manual}
      \addplot[fill=cazul!75, draw=cazul!80!black] coordinates {(Meningioma,0.8901) (Glioma,0.6259) (Hipofisario,0.8129) (Media,0.7763)};
      \addlegendentry{AG (A)}
      \addplot[fill=cbermellon!80, draw=cbermellon!80!black] coordinates {(Meningioma,0.8935) (Glioma,0.7135) (Hipofisario,0.8457) (Media,0.8175)};
      \addlegendentry{Refinado (E)}
      \addplot[fill=cverde!70, draw=cverde!80!black] coordinates {(Meningioma,0.9221) (Glioma,0.6855) (Hipofisario,0.8672) (Media,0.8250)};
      \addlegendentry{ResNet34 (F)}
    \end{axis}
  \end{tikzpicture}
  \caption{Dice por clase y medio en el conjunto de prueba para la configuración manual de referencia y los tres modelos principales del proyecto.}
  \label{fig:prueba-final}
\end{figure}

\begin{table}[htbp]
  \centering
  \caption{Contribución de cada etapa al Dice medio en el conjunto de prueba. La mejora total respecto a la referencia es de 0.1345.}
  \label{tab:contribucion}
  \begin{tabular}{llccc}
    \toprule
    \textbf{Etapa} & \textbf{Modelo} & \textbf{Dice medio} & \textbf{Ganancia} & \textbf{Proporción} \\
    \midrule
    Configuración manual             & Referencia & 0.6905 & --          & -- \\
    Metaoptimización (AG)            & A          & 0.7763 & $+0.0858$   & 63.8\,\% \\
    Refinamiento (Fases 5 y 6)       & E          & 0.8175 & $+0.0412$   & 30.6\,\% \\
    Transferencia (Fase 7)           & F          & 0.8250 & $+0.0075$   & 5.6\,\% \\
    \bottomrule
  \end{tabular}
\end{table}

El glioma concentra la mayor parte de la mejora (de 0.447 a 0.714 con el modelo E) y sigue siendo la clase limitante. El objetivo orientativo de un Dice medio de 0.85, fijado durante el desarrollo, no se alcanzó: el mejor modelo obtiene 0.8250, y solo el 3.9\,\% de los remuestreos \emph{bootstrap} alcanzan ese valor. Meningioma y tumor hipofisario superan ya 0.85 con el modelo F; alcanzar la media habría exigido un Dice de 0.76 en el glioma, por encima del 0.714 del mejor modelo para esta clase.

\subsection{Amenazas a la validez}\label{subsec:amenazas}

Las conclusiones deben interpretarse teniendo en cuenta las siguientes limitaciones:

\begin{itemize}
  \item \textbf{Una sola semilla.} Cada configuración se entrenó una única vez, y el propio AG se ejecutó una sola vez con semilla 42. Otra semilla podría haber conducido a una solución distinta, y no se dispone de la varianza entre entrenamientos.
  \item \textbf{Presupuestos desiguales.} La referencia se entrenó 40 épocas y la configuración del AG 100. La comparación con igual presupuesto en validación respalda que la mayor parte de la ganancia procede de los hiperparámetros, pero la comparación en prueba no está equilibrada.
  \item \textbf{Reutilización del conjunto de prueba.} Se usó en tres comparaciones sucesivas, por lo que el resultado final conserva cierto optimismo de selección.
  \item \textbf{Resolución del conjunto de validación.} Con 468 cortes, no permite ordenar modelos con una resolución de $\pm0.01$--$0.02$, como muestran las inversiones de orden entre validación y prueba. Esta limitación afecta tanto a la selección de puntos de control como a la aptitud del AG.
  \item \textbf{Aptitud de baja fidelidad y presupuesto reducido.} Con 46 evaluaciones y una población de 10 individuos, la exploración es limitada y existe riesgo de convergencia prematura; el sesgo de horizonte corto se ha documentado en la \cref{subsec:sesgo}.
  \item \textbf{Ausencia de una referencia algorítmica.} No se dispone de los resultados de la búsqueda aleatoria con igual presupuesto, por lo que la superioridad del AG solo está demostrada frente a la configuración manual.
  \item \textbf{Alcance de los datos.} Se trabaja con cortes bidimensionales de una única secuencia y un único conjunto público. No se ha evaluado la generalización a otros equipos u hospitales, y el \emph{bootstrap} por paciente no captura la variabilidad debida al entrenamiento.
\end{itemize}

\subsection{Trabajo futuro}\label{subsec:futuro}

Las líneas de trabajo que se derivan directamente de este estudio son las siguientes:

\begin{enumerate}
  \item \textbf{Referencias y repeticiones.} Completar la búsqueda aleatoria con igual presupuesto, repetir el AG con varias semillas y compararlo con otras metaheurísticas como la optimización por enjambre de partículas~\cite{kennedy1995pso}.
  \item \textbf{Aptitud multifidelidad.} Evaluar todos los individuos con pocas épocas y prolongar el entrenamiento solo de los más prometedores, al estilo de \emph{successive halving} o Hyperband~\cite{li2018hyperband}, para mitigar el sesgo de horizonte corto sin multiplicar el coste.
  \item \textbf{Cromosoma ampliado.} Incorporar genes para el aumento de datos (bloques espacial y de intensidad), la función de pérdida (tipo y pesos $\alpha$, $\beta$, $\gamma_T$) y la programación de la tasa de aprendizaje. En este trabajo esas dimensiones se estudiaron mediante diseño factorial; el AG podría explorarlas conjuntamente y capturar interacciones como la observada entre los pesos de la pérdida y la capacidad.
  \item \textbf{Optimización multiobjetivo.} Buscar el frente de Pareto entre precisión, número de parámetros y tiempo de entrenamiento con algoritmos como NSGA-II~\cite{deb2002nsga2}.
  \item \textbf{Evaluación más robusta.} Sustituir la partición fija por validación cruzada por paciente, para reducir el ruido de medida que hoy limita tanto la selección de modelos como la aptitud.
  \item \textbf{Glioma.} Combinar el codificador preentrenado con una pérdida más orientada a la sensibilidad ($\alpha\geq0.7$) y añadir contexto entre cortes adyacentes (2.5D o 3D).
\end{enumerate}

\subsection{Conclusiones}\label{subsec:conclusiones}

Este trabajo ha abordado la segmentación multiclase de tumores cerebrales en RM como un problema de optimización en dos niveles, y ha dado el mismo peso al entrenamiento de la red y a la metaoptimización de sus hiperparámetros. Sus conclusiones principales son las siguientes:

\begin{enumerate}
  \item \textbf{El entrenamiento de la U-Net es un problema de optimización cuyo diseño puede razonarse analíticamente.} La derivación del gradiente de la pérdida Focal Tversky (\cref{eq:gradiente-ti}) predice que los pesos $\alpha$ y $\beta$ controlan el equilibrio entre sensibilidad y precisión, y los experimentos lo confirman: la razón FN/FP del glioma pasó de 2.25 a 1.44 con $\alpha=0.7$, cayó a 0.64 al añadir aumento espacial y se equilibró en 1.11 con $\alpha=0.6$ y 32 filtros.
  \item \textbf{El algoritmo genético es la palanca más eficaz del proyecto.} Con un cromosoma real normalizado y 46 evaluaciones de bajo coste, encontró una configuración que mejora la aptitud un 15.3\,\%, eleva el Dice medio en prueba de 0.6905 a 0.7760 y converge 3.2 veces más rápido que la configuración manual. Explica casi dos tercios de la mejora total.
  \item \textbf{Una solución heurística debe analizarse críticamente.} La aptitud de tres épocas introduce un sesgo de horizonte corto que favorece redes pequeñas y es ciega a la programación del paso. Identificarlo y corregirlo (32 filtros, recocido coseno y pesos de pérdida recalibrados) elevó el modelo entrenado desde cero hasta un Dice medio de 0.8175 en prueba, conservando los valores de $\eta$ y $\lambda$ hallados por el AG.
  \item \textbf{La dinámica del descenso importa tanto como los valores de los hiperparámetros.} El recocido coseno redujo unas siete veces la oscilación del Dice de validación entre épocas (de 0.022 a 0.003), lo que convirtió la selección del modelo en una decisión fiable.
  \item \textbf{La transferencia de aprendizaje es un complemento, no un sustituto.} Un codificador ResNet34 preentrenado alcanza un Dice medio estadísticamente equivalente (0.8250), converge siete veces más rápido en épocas y resuelve la confusión entre tipos tumorales, pero no mejora la delimitación del glioma.
\end{enumerate}

Como resultado aplicado, el proyecto incluye una interfaz web desarrollada con Streamlit que permite cargar un corte de RM y visualizar en tiempo real la máscara predicha por los modelos E y F superpuesta a la imagen, con carácter demostrativo y sin finalidad clínica.


\clearpage
\printbibliography[heading=bibintoc,title={Referencias}]

\appendix
\section{Configuración de todas las ejecuciones}\label{app:configuraciones}

La \cref{tab:todas} recoge la configuración y el resultado de todas las ejecuciones citadas. Salvo indicación contraria, todas comparten profundidad $D=4$, optimizador Adam, lotes de 8 imágenes, entrada de $256\times256$ píxeles, normalización min-max, precisión mixta y semilla 42. Todas las ejecuciones posteriores a la Fase 3 usan los hiperparámetros del algoritmo genético: $\eta=3.15144\cdot10^{-4}$ y $\lambda=1.90823\cdot10^{-6}$.

\begin{table}[htbp]
  \centering
  \caption{Configuración y resultados de todas las ejecuciones. DF: Dice-Focal ($w_r=w_p=0.5$, $\gamma=2$). FT: Focal Tversky-Focal ($w_r=w_p=0.5$, $\gamma_T=4/3$) con los pesos $\alpha/\beta$ indicados. PT: parada temprana con paciencia de 20 épocas. Los aumentos se definen en la \cref{subsec:aumento}.}
  \label{tab:todas}
  \small
  \setlength{\tabcolsep}{3.4pt}
  \begin{tabular}{llcclcccc}
    \toprule
    \textbf{Ejec.} & \textbf{Red} & \textbf{Prog.} & \textbf{Épocas} & \textbf{Pérdida} & \textbf{Aumento} & \textbf{Mejor ép.} & \textbf{Dice val.} & \textbf{Dice prueba} \\
    \midrule
    Ref.$^{\dagger}$ & U-Net 16 & const. & 40 & DF & básico & 35 & 0.6538 & 0.6905 \\
    A     & U-Net 16 & const. & 100 & DF & básico & 64 & 0.7702 & 0.7760 \\
    B     & U-Net 16 & const. & 100 & DF & avanzado & 41 & 0.7595 & -- \\
    C     & U-Net 16 & const. & 100 & FT 0.7/0.3 & básico & 52 & 0.7672 & -- \\
    D     & U-Net 16 & const. & 100 & FT 0.7/0.3 & avanzado & 97 & 0.7675 & -- \\
    D-ext & U-Net 16 & const. & 42 (PT) & FT 0.7/0.3 & avanzado & 22 & 0.7440 & -- \\
    S     & U-Net 16 & const. & 71 (PT) & FT 0.7/0.3 & esp. avanz. & 51 & 0.7751 & -- \\
    I     & U-Net 16 & const. & 64 (PT) & FT 0.7/0.3 & intensidad & 44 & 0.7454 & -- \\
    E16   & U-Net 16 & coseno & 150 & FT 0.6/0.4 & esp. avanz. & 95 & 0.7693 & -- \\
    E     & U-Net 32 & coseno & 150 & FT 0.6/0.4 & esp. avanz. & 72 & 0.7612 & 0.8175 \\
    F$^{\ast}$ & ResNet34 & coseno & 100 & FT 0.6/0.4 & esp. avanz. & 84 & 0.8221 & 0.8250 \\
    \bottomrule
  \end{tabular}

  \smallskip
  \parbox{0.92\linewidth}{\footnotesize U-Net 16 y U-Net 32: U-Net propia con 16 o 32 filtros base. $^{\dagger}$La referencia usa $\eta=10^{-3}$ y $\lambda=10^{-4}$; el resto, los valores del AG. $^{\ast}$U-Net de Segmentation Models PyTorch con codificador ResNet34 preentrenado, cuya tasa de aprendizaje es diez veces menor que la del decodificador.}
\end{table}

\section{Reproducibilidad}\label{app:reproducibilidad}

El código fuente del proyecto está organizado en módulos de datos (\texttt{src/dataset}), modelos y pérdidas (\texttt{src/models}), optimización (\texttt{src/optimization}) y utilidades de entrenamiento y evaluación (\texttt{src/utils}), con guiones ejecutables en \texttt{scripts/} y la interfaz en \texttt{app/}. El algoritmo genético se implementa en \texttt{src/optimization/genetic.py} y su función de aptitud en \texttt{src/optimization/fitness.py}. Las órdenes siguientes reproducen las ejecuciones principales desde la raíz del repositorio:

\begin{small}
\begin{verbatim}
python -m scripts.run_optimization --epochs 3 --device auto --amp

python -m scripts.train_baseline --epochs 100 --learning-rate 0.000315144 \
    --weight-decay 1.90823e-06 --base-filters 16 --depth 4 --batch-size 8 \
    --seed 42 --device auto --amp

python -m scripts.train_baseline --epochs 150 --learning-rate 0.000315144 \
    --weight-decay 1.90823e-06 --base-filters 32 --depth 4 --batch-size 8 \
    --seed 42 --augmentation advanced_spatial --loss focal_tversky_focal \
    --tversky-alpha 0.6 --tversky-beta 0.4 --lr-schedule cosine \
    --device auto --amp

python -m scripts.train_baseline --model smp_unet --encoder resnet34 \
    --encoder-weights imagenet --encoder-lr-multiplier 0.1 --epochs 100 \
    --learning-rate 0.000315144 --weight-decay 1.90823e-06 --batch-size 8 \
    --seed 42 --augmentation advanced_spatial --loss focal_tversky_focal \
    --tversky-alpha 0.6 --tversky-beta 0.4 --lr-schedule cosine \
    --device auto --amp

python -m scripts.evaluate_test_set
streamlit run app/main.py
\end{verbatim}
\end{small}

Las órdenes corresponden, por este orden, a la búsqueda del algoritmo genético, al reentrenamiento de la configuración del AG (ejecución A), al modelo refinado (E), al modelo con transferencia de aprendizaje (F), a la evaluación en el conjunto de prueba y a la interfaz de inferencia. Los datos de todas las figuras de esta memoria se encuentran en la carpeta \texttt{data/} del propio documento.


\end{document}
